% Self-contained source: the bibliography is included at the end of this file.
% Compile twice with:
%   pdflatex formal_categories_and_distributors.tex
% Alternatively use: latexmk -pdf formal_categories_and_distributors.tex
\documentclass[11pt,reqno]{amsart}
\usepackage[T1]{fontenc}
\usepackage{lmodern}
\usepackage{amsmath,amssymb,amsthm,mathtools}
\usepackage{xcolor}
\usepackage{tikz-cd}
\usepackage{booktabs,array}
\usepackage[margin=1.12in]{geometry}
\usepackage{microtype}
\usepackage{enumitem,etoolbox}
% Keep proof steps compact; statement lists retain their original formatting.
\AtBeginEnvironment{proof}{%
  \emergencystretch=1em
  \setlist[enumerate,1]{label=(\arabic*),leftmargin=1.85em,labelsep=.45em,
    itemsep=.28\baselineskip plus .1\baselineskip minus .05\baselineskip,
    topsep=.35\baselineskip,parsep=0pt,partopsep=0pt,beginpenalty=10000}%
}
\usepackage[colorlinks=true,linkcolor=blue!55!black,citecolor=blue!55!black,urlcolor=blue!55!black]{hyperref}
\usepackage[nameinlink,noabbrev]{cleveref}
\usetikzlibrary{arrows.meta,quotes}
\tikzcdset{row sep=2.2em,column sep=2.8em}
\tikzset{proarrow/.style={"{\scriptscriptstyle /}" marking}}
\theoremstyle{plain}
\newtheorem{theorem}{Theorem}[section]
\newtheorem{proposition}[theorem]{Proposition}
\newtheorem{lemma}[theorem]{Lemma}
\newtheorem{corollary}[theorem]{Corollary}
\theoremstyle{definition}
\newtheorem{definition}[theorem]{Definition}
\newtheorem{example}[theorem]{Example}
\theoremstyle{remark}
\newtheorem{remark}[theorem]{Remark}
\numberwithin{equation}{section}
\newcommand{\D}{\mathbb D}
\newcommand{\Pcal}{\mathbb P}
\newcommand{\Bcal}{\mathbb B}
\newcommand{\V}{\mathcal V}
\newcommand{\E}{\mathsf E}
\newcommand{\Q}{\mathsf Q}
\newcommand{\Ccal}{\mathbb C}
\newcommand{\Kcal}{\mathbb K}
\newcommand{\Ddist}{\mathbb D_{\mathrm{dist}}}
\newcommand{\FDist}{\mathsf{Dist}_{\mathrm{poly}}}
\newcommand{\Up}{U_{\mathrm{poly}}}
\newcommand{\harr}{\mathrel{\nrightarrow}}
\newcommand{\cell}{\mathrel{\Longrightarrow}}
\newcommand{\emp}[1]{\varnothing_{#1}}
\newcommand{\id}{\mathrm{id}}
\newcommand{\op}{\mathrm{op}}
\newcommand{\vop}{\mathrm v}
\newcommand{\Mod}{\operatorname{Mod}_{\mathrm{CS}}}
\newcommand{\Mon}{\operatorname{Mon}}
\newcommand{\FCat}{\operatorname{FCat}}
\newcommand{\Bicomod}{\operatorname{Bicomod}}
\newcommand{\Comon}{\operatorname{Comon}}
\newcommand{\Hom}{\operatorname{Hom}}
\newcommand{\Set}{\mathbf{Set}}
\newcommand{\Ab}{\mathbf{Ab}}
\newcommand{\Z}{\mathbb Z}
\newcommand{\Icat}{\mathsf 1}
\newcommand{\Forest}{\mathsf{For}}
\newcommand{\Raw}{\mathsf{Path}}
\newcommand{\eps}{\varepsilon}
\newcommand{\wh}{\mathbin{\odot}}
\newcommand{\bal}{\operatorname{Bal}}
\DeclareMathOperator{\Eq}{Eq}
\title[Formal categories and distributors]{Non-Cartesian Internal Categories\\and Distributors without Regularity:\\A Poly-Virtual Approach}
\date{}
\subjclass[2020]{Primary 18N10; Secondary 18M05, 18M65, 18M85, 16T15.}
\keywords{internal category, formal distributor, bicomodule, cotensor, virtual equipment, polycategory, localization}

\begin{document}
\begin{abstract}
Aguiar's non-Cartesian internal categories use bicomodule cotensors for
composition. We define formal internal categories in every monoidal
category on actual bicomodule carriers, using a poly-virtual calculus
with finite paths on both boundaries. A forest envelope and localization
make intrinsic globular contextually universal comparisons reversible.
Substitution presheaves and backward evaluation prove conservativity,
including equivariant maps over changing comonoid bases. Stable cotensors
recover the usual category structures, functors, and bimodule actions.

Formal categories then form the objects of an equipment of distributors
on bicomodule-path carriers. Cartesian factorization constructs their
restrictions; a second conservative completion supplies many-sided cells
and coherent boundary transport. Contextually preserved action
coequalizers represent distributor composition. Finite group-like bases
in abelian groups recover enriched categories and distributors despite
ambient nonregularity. A stronger triangular example has a nonzero action
on one actual bicomodule whose input path has no contextual atomic
representative: restriction along a formal functor preserves the action
while destroying stability of its cotensor. The general completion also
preserves virtual monoids and their generalized-multicategory examples.
\end{abstract}
\maketitle
\tableofcontents

\section{Introduction}\label{sec:intro}

An internal category has an object of objects, an object of arrows, and
composition on composable pairs. In a monoidal category, Aguiar expresses
these data by a comonoid $C$, a $C$--$C$ bicomodule $A$, and maps
\[
 m:A\square_C A\longrightarrow A,\qquad u:C\longrightarrow A,
\]
with associative and unital multiplication
\cite[Definition 2.3.1]{Aguiar1997}. The comonoid and its coactions replace
the object of objects and the source and target maps; the cotensor replaces
the pullback of composable arrows. This gives a formulation that also
makes sense when the monoidal product is not Cartesian.

Cotensor composition requires equalizers and their compatibility with
tensoring. Aguiar's regularity condition supplies this compatibility for
every equalizer in the ambient monoidal category
\cite[Definition 2.1.1]{Aguiar1997}. However, even
$(\Ab,\otimes_{\mathbb Z},\mathbb Z)$ fails this condition. As
\cref{prop:Ab} shows, the finite group-like comonoids in $\Ab$
nevertheless admit stable bicomodule composition and encode the expected
finite-object enriched categories. A general definition should therefore
separate the specification of category operations from the existence of
all ambient horizontal composites.

Cruttwell and Shulman's virtual framework provides a precedent for this
separation. It treats enriched categories, internal categories, and
generalized multicategories through cells with several horizontal inputs
and one output \cite{CruttwellShulman2010}. In particular, its virtual
equipment of enriched profunctors is available without the coends used
to represent horizontal profunctor composition. The relevant analogy is
between these coends and bicomodule cotensors. Coends here compose
profunctors between enriched categories; the composition operations inside
an enriched category are already part of its defining data.

The guiding question of this paper is therefore the following: can
non-Cartesian internal categories be defined over an arbitrary monoidal
category, on actual bicomodule carriers, so that stable cotensors recover
Aguiar's multiplication and changing-base functors, and can these formal
categories themselves be the objects of an equipment of distributors?
The primitive maps
available without cotensors have the form
\[
 N\longrightarrow M_1\otimes\cdots\otimes M_n,
\]
subject to balance and equivariance equations. They describe cones into a
prospective composite. Their arity is opposite to that of multiplication.
Consequently, the existing virtual bicomodule calculus alone does not yet
supply the desired forward category operations.

We construct a two-sided cell calculus in which both boundaries are finite
paths of bicomodules. A formal internal category has an endobicomodule
$A:C\harr C$ and cells
\begin{equation}\label{eq:shape}
 \mu:(A,A)\cell A,\qquad \eta:\emp C\cell A
\end{equation}
satisfying the monad equations. The construction makes each existing intrinsic globular
contextually universal composite comparison reversible. It thereby allows
represented cotensors to be used on either boundary, while paths without
representatives remain legitimate boundaries. Its cells are specified by
finite expressions and relations, so the definition does not depend on
the existence of the object $A\square_C A$.

\begin{theorem}[Internal categories without regularity]\label{thm:main}
Every monoidal category $\V$ determines a poly-virtual equipment $\Q(\V)$
and a category $\FCat(\Q(\V))$ with the following properties.
\begin{enumerate}
\item The objects of $\Q(\V)$ are comonoids, its vertical arrows are
comonoid maps, and its designated atomic horizontal arrows are actual
bicomodules. Its unary squares between designated atomic bicomodules are
exactly equivariant bicomodule maps.
More generally, all primitive balanced cells are preserved and reflected.
\item A formal internal category consists of an atomic endobicomodule
and cells as in \cref{eq:shape}, with associativity and unit equations
given by tree substitution. Homomorphisms have ordinary forward variance
and may change the comonoid base.
\item On bicomodule paths represented by cotensor cones preserved under
tensoring, these structures, their homomorphisms, and their bimodule
actions recover the corresponding cotensor structures and equations.
In particular, for regular $\V$ the category $\FCat(\Q(\V))$ is
isomorphic to Aguiar's category of internal categories and internal functors.
\end{enumerate}
\end{theorem}

The construction applies more generally to virtual equipments. Here
\textbf{contextually universal} means that a comparison represents its
source path in every compatible horizontal context, including contexts
that pass through different objects. The technical result is a conservative
completion that recognizes all such comparisons intrinsically.

\begin{theorem}[Conservative completion]\label{thm:maincompletion}
Every virtual equipment $\D$ has a merge realization $\E(\D)$ and
a poly-virtual reduct $\Up\E(\D)$ on the same objects, vertical
arrows, and designated atomic horizontal arrows, with natural bijections
\[
 \D_{f,g}(\Gamma;H)\cong\E(\D)_{f,g}(\Gamma;H)
\]
for all original paths $\Gamma$, atomic outputs $H$, and compatible
vertical boundaries. The realization makes every intrinsic globular
contextually universal comparison, and hence every horizontal contextual
instance of one, reversible. It also admits coherent transport of horizontal
boundaries. The reduct retains all cells and equalities and the tree
operations; it preserves and reflects formal categories and their
homomorphisms on atomic carriers.
\end{theorem}

The bicomodule instance is
\[
 \Q(\V)=\Up\bigl(\E(\Mod(\V^\op))^{\vop}\bigr).
\]
Here $\Mod$ is Leinster's monoids-and-bimodules construction
\cite[Section 5.3]{Leinster2004}, in the virtual-double-category formulation
of Cruttwell and Shulman \cite[Definition 2.8]{CruttwellShulman2010}. It
supplies the initial balanced calculus; the completion
and the reversal of squares supply its two-sided, forward formulation.

The distributor construction uses the same completion a second time.
Let $\mathbb B(\V)=\E(\Ccal(\V))^{\vop}$ be the merge realization,
and regard its finite bicomodule paths as individual horizontal arrows
of an underlying virtual equipment $\Kcal(\V)$. We form
$\Mod(\Kcal(\V))$ and retain exactly the monoids whose carriers are
original atomic bicomodules. The restriction here is on objects: the
horizontal carriers of the resulting equipment $\Ddist(\V)$ remain
finite bicomodule paths. This makes the required distributor restrictions
available while retaining all original atomic bimodules. Define
\[
 \FDist(\V)=\Up\E(\Ddist(\V)).
\]

\begin{theorem}[Formal distributors without regularity]\label{thm:maindistributors}
For every monoidal category $\V$, $\FDist(\V)$ is a poly-virtual
equipment with the following properties.
\begin{enumerate}
\item Its objects are the formal internal categories in $\Q(\V)$,
and its vertical category is exactly $\FCat(\Q(\V))$.
\item Its designated horizontal arrows are distributors with associative,
unital, commuting actions on finite bicomodule-path carriers. On a carrier
consisting of one actual bicomodule, these are exactly the atomic bimodule
actions in $\Q(\V)$.
\item Balanced single-output distributor cells are preserved and
reflected by the completion, including all their changing-base instances.
In particular, unary cells between distributors carried by one actual
bicomodule are exactly equivariant bicomodule maps preserving both actions.
\item Formal functors have companions and conjoints, yielding coherent
transport of distributor boundaries. Every intrinsic globular contextually
universal distributor-composition comparison becomes reversible.
\end{enumerate}
\end{theorem}

There are two distinct representation problems. Cotensors represent
paths of bicomodules over comonoid bases. Relative tensor products of
distributors represent balancing over the actions of an intervening
formal category. We give a contextual coequalizer construction for the
latter and verify it in the finite group-like $\Ab$ sector. Neither
construction requires a missing representative to be an object of $\V$.

The distinction occurs inside an action in \cref{strong:thm:strong}.
There, a formal category with carrier $A\cong\Z^5$ acts on an actual
bicomodule $M\cong(\Z/4)^2$, although $(A,M)$ has no contextual atomic
representative. Both categories used to construct the action have stable
finite composition paths. Restriction along a formal functor produces
the action from a stable one, showing that stability of action cotensors
need not survive restriction, even over an identity comonoid map.

The term \textbf{poly-virtual equipment} has the precise meaning given in
\cref{def:poly,def:equipment}. Its equipment structure transports
boundaries, and a transported boundary can itself be a path. The tree and
merge operations are related to those of \cite{Hadzihasanovic2019}; empty
boundaries and a vertical category are included here. The relation with
implicit structures is discussed in \cref{sec:comparison}.

The construction and conservativity occupy
\cref{sec:bicomod,sec:typed,sec:monads,sec:forest,sec:localization,sec:global,sec:completion};
\cref{sec:formalinternal,sec:recovery,sec:changing} give the bicomodule
application and cotensor recovery. Distributors and their representations
are developed in
\cref{sec:distributors,sec:distributorcompletion,sec:distributorrecovery}.
The examples in \cref{sec:examples,sec:unrepresented} separate ambient
nonregularity from failure of representation at an action boundary.
Further applications and comparison with other envelopes follow in
\cref{sec:kleisli,sec:comparison}.

\subsection*{Conventions}

Horizontal paths are written in traversal order: $HK$ means first $H$,
then $K$. Composition of ordinary maps and of cells is written in functional
order: $ba=b\circ a$. Horizontal juxtaposition of cells is $a\wh b$.
The path of length zero at $X$ is $\emp X$; a represented unit is denoted
$U_X$. A semicolon separates the two boundaries of a cell.
Associators and unitors in a monoidal category are inserted by their
canonical coherence maps; no braiding is used unless explicitly stated.
The notation $\V^\op$ reverses maps and leaves the order of tensor factors
unchanged. Size is discussed in \cref{app:size}.
We write $\mathsf A=(X,A,\mu_A,\eta_A)$ for a formal category and
$\mathsf F=(f,F)$ for a formal functor, separating the whole structure
from its carrier and its underlying square. The word \textbf{atomic} is
relative to the designated horizontal graph: an atomic arrow of the
distributor equipment is one distributor, whose carrier in the original
bicomodule graph can be a path.
We use the slashed proarrow $\harr$ of
\cite[Definition 2.1]{CruttwellShulman2010}; its slash is a proarrow
marker. Double-category diagrams use slashed horizontal arrows and a
centered double arrow for the cell. In diagrams whose vertices are paths,
ordinary arrows denote cells in the merge realization. Commutativity
asserts equality of composites; the category and action laws compare
tree composites. Dashed arrows denote unique factorizations, except
where explicitly labelled $\nexists$.

\subsection*{Note on use of AI}
Generative AI systems were used during the research and development of this manuscript. Responsibility for the mathematical content,
arguments, and conclusions rests with the author.


\section{Bicomodules and the missing composition}\label{sec:bicomod}

\subsection{Virtual composition and its enriched precedent}
A virtual double category has objects, vertical maps, horizontal arrows,
and cells from a finite path to one horizontal arrow. Its composition is
substitution of compatible cells. A cell is globular when its exterior
vertical maps are identities. A composite of a horizontal path, when
it exists, represents these substitutions in every surrounding path
context. A virtual equipment has represented horizontal units and
cartesian restrictions along pairs of vertical arrows
\cite[Sections 2, 5, and 7]{CruttwellShulman2010}.

For enriched profunctors $P:\mathcal A\harr\mathcal B$ and
$R:\mathcal B\harr\mathcal C$, the represented composite has
components
\[
 (P\wh R)(a,c)=\int^{b\in\mathcal B}P(a,b)\otimes R(b,c).
\]
Balanced families of maps out of these tensor factors define the virtual
cells before the coend is formed. This is the distinction between the
virtual equipment and its represented composites
\cite[Examples 5.6 and 7.7]{CruttwellShulman2010}.
For bicomodules the basic balanced maps point into tensor products.
We now specify those maps without taking their equalizers.

\subsection{Comonoids, coactions, and balanced cells}

\begin{definition}
A comonoid in $\V$ is an object $C$ with coassociative
$\delta_C:C\to C\otimes C$ and counit $\eps_C:C\to I$.
A comonoid map preserves both maps. A bicomodule $M:C\harr D$ has
coactions
\[
 \lambda_M:M\to C\otimes M,\qquad
 \rho_M:M\to M\otimes D
\]
which are counital and coassociative and commute:
\begin{equation}\label{eq:bicomodcommute}
\begin{tikzcd}[column sep=3.4em]
M \arrow[r,"\lambda_M"] \arrow[d,"\rho_M"'] &
C\otimes M \arrow[d,"1_C\otimes\rho_M"]\\
M\otimes D \arrow[r,"\lambda_M\otimes1_D"'] & C\otimes M\otimes D.
\end{tikzcd}
\end{equation}
A map $F:M\to N$ over comonoid maps $f:C\to C'$, $g:D\to D'$
is equivariant if
\begin{equation}\label{eq:equivariant}
\begin{tikzcd}[column sep=2.4em]
M \arrow[r,"F"] \arrow[d,"\lambda_M"'] & N \arrow[d,"\lambda_N"]\\
C\otimes M \arrow[r,"f\otimes F"'] & C'\otimes N
\end{tikzcd}
\qquad
\begin{tikzcd}[column sep=2.4em]
M \arrow[r,"F"] \arrow[d,"\rho_M"'] & N \arrow[d,"\rho_N"]\\
M\otimes D \arrow[r,"F\otimes g"'] & N\otimes D'.
\end{tikzcd}
\end{equation}
\end{definition}

Our horizontal arrow points from the left coaction base to the right
coaction base. References using the opposite horizontal direction write
composite bicomodules in the opposite order. This convention fixes the
order of all tensor factors below.

For $M:C\harr D$ and $N:D\harr E$, the prospective cotensor
is the equalizer of
\[
 \rho_M\otimes1_N,\;1_M\otimes\lambda_N:
 M\otimes N\rightrightarrows M\otimes D\otimes N.
\]
Even if this equalizer is unavailable, maps equalizing the displayed pair
can still be specified. Longer paths impose one such balance equation
at each interior base, together with the exterior coaction equations.

\begin{definition}[Balanced cells]\label{def:balanced}
Let $M_i:C_{i-1}\harr C_i$ for $1\leq i\leq n$, let
$N:D\harr E$, and let $a:D\to C_0$, $b:E\to C_n$ be comonoid
maps. For $n>0$, set $T=M_1\otimes\cdots\otimes M_n$. A balanced
cell is a map $\phi:N\to T$ with exterior equivariance
\begin{equation}\label{eq:exterior}
\begin{tikzcd}[column sep=2.3em]
N \arrow[r,"\phi"] \arrow[d,"\lambda_N"'] &
T \arrow[d,"\lambda_{M_1}\otimes1"]\\
D\otimes N \arrow[r,"a\otimes\phi"'] & C_0\otimes T
\end{tikzcd}
\qquad
\begin{tikzcd}[column sep=2.3em]
N \arrow[r,"\phi"] \arrow[d,"\rho_N"'] &
T \arrow[d,"1\otimes\rho_{M_n}"]\\
N\otimes E \arrow[r,"\phi\otimes b"'] & T\otimes C_n
\end{tikzcd}
\end{equation}
and one interior balance at each $C_i$. Write $T_i$ for $T$ with
$C_i$ inserted between $M_i$ and $M_{i+1}$, and let
$r_i=1\otimes\rho_{M_i}\otimes1$ and
$\ell_i=1\otimes\lambda_{M_{i+1}}\otimes1$. Balance says that
the two routes have a common composite $\psi_i$:
\begin{equation}\label{eq:balance}
\begin{tikzcd}[column sep=3.4em]
N \arrow[r,"\phi"] \arrow[dr,"\psi_i"'] &
T \arrow[d,shift left=.7ex,"r_i"]
  \arrow[d,shift right=.7ex,"\ell_i"']\\
& T_i.
\end{tikzcd}
\qquad (1\leq i<n).
\end{equation}
For $n=0$, the path is $\emp C$, $a:D\to C$, $b:E\to C$, and
a balanced cell is a map $\phi:N\to I$ making
\begin{equation}\label{eq:nullbalance}
\begin{tikzcd}[column sep=3em]
N \arrow[r,"\rho_N"] \arrow[d,"\lambda_N"'] &
N\otimes E \arrow[d,"\phi\otimes b"]\\
D\otimes N \arrow[r,"a\otimes\phi"'] & C
\end{tikzcd}
\end{equation}
commute, with unitors understood.
\end{definition}

\subsection{The virtual equipment already available}
The balanced maps are closed under substitution. Their orientation as
single-output cells is forced by that substitution law: the map from
$N$ into a tensor of $M_i$ is regarded as a cell with output $N$.
This produces the following instance of the monoids-and-bimodules
construction of Leinster, in the virtual-double-category formulation of
Cruttwell and Shulman \cite[Section 5.3]{Leinster2004}
\cite[Definition 2.8]{CruttwellShulman2010}.

\begin{proposition}\label{prop:modinput}
Comonoids, opposite comonoid maps, bicomodules, and the balanced cells
of \cref{def:balanced}, regarded as cells
$(M_1,\ldots,M_n)\cell N$, form the virtual equipment
\[
 \Ccal(\V)=\Mod(\V^\op).
\]
This construction exists for every monoidal $\V$.
\end{proposition}
The balanced map $\phi:N\to M_1\otimes\cdots\otimes M_n$ has profile
\[
\begin{tikzcd}[column sep=6em,row sep=3.4em]
C_0 \arrow[r,proarrow,"{(M_1,\ldots,M_n)}",""{name=top,below}]
  \arrow[d,"a^{\op}"'] & C_n \arrow[d,"b^{\op}"]\\
D \arrow[r,proarrow,"N"',""{name=bottom,above}] & E
\arrow[Rightarrow,from=top,to=bottom,shorten <=2pt,shorten >=2pt,"\phi"]
\end{tikzcd}
\]
in the opposite comonoid category.
\begin{proof}\leavevmode
\begin{enumerate}
\item \textbf{Identify the module construction.} Regard $\V^\op$ as the one-object virtual double category whose cells
are maps into tensor products in $\V$. The construction $\Mod$ requires no
additional exactness condition on this virtual double category
\cite[Section 5.3]{Leinster2004}\cite[Remark following Definition 2.8]{CruttwellShulman2010}. A monoid there is precisely a
comonoid in $\V$, and a bimodule is precisely a bicomodule. The exterior
action equations and interior balancing equations in the module
construction become \cref{eq:exterior,eq:balance}; the zero-arity
action equation becomes \cref{eq:nullbalance}.

\item \textbf{Check substitution at nonempty blocks.} Explicitly, substitution composes $\phi$ with the tensor product of the
maps substituted into its factors. An interior equation inside a block
is the corresponding equation of the substituted map, tensored with
the other maps. At an interface between two nonempty blocks,
\cref{eq:exterior} moves the coactions through those maps and the
remaining equality is \cref{eq:balance} for $\phi$.

\item \textbf{Handle empty blocks.} An empty block
is removed using its map to $I$; equation~\eqref{eq:nullbalance}
identifies the two induced maps to its vertex comonoid, so the same
interface equality remains true. Repeating this for successive empty
blocks proves the statement also when every block is empty.

\item \textbf{Verify the exterior and substitution laws.} The exterior
equations follow by the first and last such calculations. Associativity
and identities now follow from tensor functoriality and composition in
$\V$.

\item \textbf{Construct restrictions.} Restrictions are explicit. To restrict $N:D\harr E$ along vertical
arrows $C\to D$ and $C'\to E$ in the opposite comonoid category,
write their actual maps as $a:D\to C$, $b:E\to C'$. Retain the
object $N$ and replace its coactions by
$(a\otimes1)\lambda_N$ and $(1\otimes b)\rho_N$.

\item \textbf{Verify cartesianness.} The identity map on $N$ is the restriction cell. The equations for
factorization through it are literally \cref{eq:exterior} with these
coactions, proving cartesianness.

\item \textbf{Supply units.} Units are proved in
\cref{lem:bicomodunit} below. This also directly verifies the units
and restrictions supplied abstractly by
\cite[Propositions 5.5 and 7.4]{CruttwellShulman2010}.\qedhere
\end{enumerate}
\end{proof}

\begin{lemma}[The nullary representative]\label{lem:bicomodunit}
For a comonoid $C$, its regular bicomodule
$(C,\delta_C,\delta_C)$ represents $\emp C$ in $\Ccal(\V)$.
The universal cell $\emp C\cell C$ is represented by $\eps_C$.
\end{lemma}
\begin{proof}\leavevmode
\begin{enumerate}
\item \textbf{Construct the nullary factor.} First consider a globular nullary cell $\phi:N\to I$. By
\cref{eq:nullbalance},
\[
 q=(1_C\otimes\phi)\lambda_N
   =(\phi\otimes1_C)\rho_N:N\to C.
\]

\item \textbf{Verify equivariance.} Coassociativity of $\lambda_N$ gives
$\delta_Cq=(1_C\otimes q)\lambda_N$, and that of $\rho_N$ gives
$\delta_Cq=(q\otimes1_C)\rho_N$; in each equation use the
corresponding expression for $q$. Thus $q$ is a bicomodule map.

\item \textbf{Prove nullary uniqueness.} Counitality gives the factorization triangle
\[
\begin{tikzcd}[column sep=3em]
N \arrow[r,dashed,"\exists!\,q"] \arrow[dr,"\phi"'] &
C \arrow[d,"\eps_C"]\\
& I.
\end{tikzcd}
\]
Conversely a bicomodule map
$q:N\to C$ satisfies both displayed expressions with
$\phi=\eps_Cq$, proving uniqueness.

\item \textbf{Insert the unit in a context.} For a contextual instance, take a balanced map
$\phi:N\to\bigotimes L\otimes\bigotimes R$ in which the two
paths meet at $C$. If $L$ is nonempty, insert $C$ by the right
coaction of its last factor. If $R$ is nonempty, one can instead use
the left coaction of its first factor. Interior balance says these
two constructions agree when both apply. If both are empty, use the
map $q$ just constructed. If only one side is present, the exterior
equation supplies the missing endpoint calculation.

\item \textbf{Check the contextual equations.} Coassociativity
gives balance at the two new interfaces, and commuting coactions give
the unchanged exterior equations.

\item \textbf{Prove contextual uniqueness.} Removing the inserted factor by
$\eps_C$ recovers $\phi$. Conversely, either balance equation
adjacent to an inserted $C$, followed by its counit, forces this
insertion formula. Thus the factorization is unique in all contexts.

\item \textbf{Allow changing exterior maps.} For nonidentity exterior maps, corestrict the two exterior coactions
by the explicit restriction construction above. The resulting
factorization is the globular one just proved, with the same underlying
maps in $\V$.\qedhere
\end{enumerate}
\end{proof}

\section{Two-sided cells and tree pastings}\label{sec:typed}

\begin{definition}[Boundary signature]
A boundary signature consists of a category $\mathcal C$ and a directed
graph of horizontal arrows on its objects. A horizontal path
$\Gamma=(H_1,\ldots,H_n):X_0\harr X_n$ includes its entire sequence of
vertices. A cell profile is $(f,g;\Gamma;\Delta)$, where
$\Gamma:X\harr Y$, $\Delta:X'\harr Y'$, $f:X\to X'$, and
$g:Y\to Y'$ lie in $\mathcal C$. We depict it as
\[
\begin{tikzcd}[column sep=5em,row sep=3.2em]
X \arrow[r,proarrow,"\Gamma",""{name=top,below}] \arrow[d,"f"'] &
Y \arrow[d,"g"]\\
X' \arrow[r,proarrow,"\Delta"',""{name=bottom,above}] & Y'
\arrow[Rightarrow,from=top,to=bottom,shorten <=2pt,shorten >=2pt,"\alpha"]
\end{tikzcd}
\]
A globular profile has $f=\id_X$ and $g=\id_Y$.
\end{definition}

We specify the pasting shapes independently of any cells that will label
them. Start with rectangles with finitely many ordered horizontal ports on
their upper and lower sides. The horizontal ports are labelled by arrows
and their vertices; the two lateral sides are labelled by composable
vertical paths. A rectangular scheme is a finite expression formed by
horizontal and vertical pasting of such rectangles, identity strips on
horizontal paths, and empty horizontal strips on vertical arrows. Pasting
requires literal agreement of the full common boundary. In particular,
horizontal pasting requires equality of the common vertical arrow, not
merely agreement of its endpoints. Empty strips compose as the arrows of
$\mathcal C$ do.

The \textbf{wire graph} of a scheme is the finite multigraph whose vertices
are its labelled rectangles and whose edges are its internal atomic horizontal
wires; structural identity strips are contracted. Acyclicity is understood for
this multigraph, so in particular two parallel edges form a cycle. A
\textbf{tree scheme} is a rectangular scheme whose wire graph is nonempty,
connected, and acyclic, and which has no separate
horizontal wire component meeting no labelled rectangle. The last
condition excludes adjoining unused through-wires as a tree operation.
For the poly-virtual calculus the permitted schemes with no labelled
rectangle are a single identity wire and an empty horizontal square on
a vertical arrow. Call these \textbf{elementary structural schemes}.
Two schemes describe the same pasting
when they differ by planar deformation, insertion or deletion of structural
identity strips, associativity of pasting, or interchange of rectangular
subdivisions. The arity of an internal seam is its number of atomic wires.
Thus two boxes joined along two wires are a merge scheme, not a tree scheme.

\begin{lemma}\label{lem:shapes}
Substitution of a tree scheme into a labelled rectangle of another tree
scheme again gives a tree scheme. Insertion of an elementary structural
scheme gives a tree scheme or an elementary structural scheme. Such
substitution is associative and has one-rectangle schemes as identities.
Evaluation orders related by contracting connected rectangular subtrees
describe the same substituted scheme.
\end{lemma}
\begin{proof}\leavevmode
\begin{enumerate}
\item \textbf{Preserve the tree shape.} Rectangular substitution preserves every boundary label and replaces one
graph vertex by a tree. Each old incident edge attaches to its designated
port. The resulting graph is connected. A cycle entirely in the inserted
graph is impossible; any other cycle would contract to a cycle in the old
graph.

\item \textbf{Prove substitution associativity.} Substitution of finite labelled rectangles is associative before
passing to planar equivalence. The equivalence respects substitution since
each of its generating moves can be performed inside the substituted
rectangle.

\item \textbf{Check structural substitutions.} A single identity wire removes a vertex having one input and
one output, preserving connectivity and acyclicity. An empty structural
square can replace only a zero-port vertex, which in a connected scheme
must be its sole vertex. These observations prove the structural cases.

\item \textbf{Obtain the evaluation identity.} The last assertion is the resulting associativity identity.\qedhere
\end{enumerate}
\end{proof}

\begin{definition}[Poly-virtual and merge calculi]\label{def:poly}
A \textbf{poly-virtual double calculus} on a boundary signature assigns a
collection $\Pcal_{f,g}(\Gamma;\Delta)$ to every profile and evaluates every
labelled tree scheme. It supplies the cells for single identity wires and
structural empty squares, and evaluates elementary structural schemes.
Evaluation respects the equivalences of
schemes and substitution of schemes into rectangles, and evaluation of a
one-rectangle scheme is its label. A map preserves this data.

A \textbf{merge calculus} has evaluation of all rectangular schemes, with
the same equations, including all structural schemes on arbitrary paths.
Its \textbf{poly-virtual reduct} $\Up\Bcal$ keeps the
same cell collections and equalities and restricts the specified evaluation
operations to tree and elementary structural schemes. In particular an
identity cell on a longer path remains a cell, but its empty multiwire
scheme is not a designated poly-virtual operation. Original horizontal arrows are
called \textbf{atomic} when a calculus is presented using their paths.
\end{definition}

This is an algebraic definition: \cref{lem:shapes} constructs the
multicategory of allowed operations, and a calculus is an algebra for those
operations. It does not require every many-sided pasting to be evaluable.

For clarity, an elementary globular cut can be described without pictures.
Let
\[
 \alpha:\Gamma\cell(\Delta_0,M,\Delta_1),\qquad
 \beta:(\Lambda_0,M,\Lambda_1)\cell\Sigma.
\]
The planar two-rectangle cut is allowed when all endpoints match and
\begin{equation}\label{eq:cutplacement}
 (\Delta_0=\varnothing\text{ or }\Lambda_0=\varnothing),\qquad
 (\Delta_1=\varnothing\text{ or }\Lambda_1=\varnothing).
\end{equation}
Its profile is
\[
 (\Lambda_0,\Gamma,\Lambda_1)\cell
 (\Delta_0,\Sigma,\Delta_1).
\]
Indeed, these conditions say that on either side of the common wire only
one box has an exposed boundary to traverse. If both corresponding
contexts are nonempty, the putative exterior is not the stated pair of
directed paths. One useful instance is
\begin{equation}\label{eq:cornercut}
 \frac{\Gamma\cell(\Delta,M)\qquad (M,\Lambda)\cell\Sigma}
 { (\Gamma,\Lambda)\cell(\Delta,\Sigma)}.
\end{equation}

\begin{proposition}\label{prop:fragments}
The single-output fragment of a poly-virtual calculus is a virtual double
category. Its single-input fragment is the vertically reversed version of
a virtual double category. Ordinary virtual substitution, including its
changing-base and nullary instances, is an allowed tree operation.
\end{proposition}
\begin{proof}\leavevmode
\begin{enumerate}
\item \textbf{Identify the substitution tree.} An outer single-output rectangle and a compatible family of single-output
rectangles above its input ports form a rooted tree. Shared lateral maps
are precisely the compatibility data of virtual substitution.

\item \textbf{Verify the virtual laws.} Replacing
vertices of this tree gives associativity by \cref{lem:shapes}; structural
strips give the identity laws. A nullary vertex has no incoming wire and
causes no change in this argument.

\item \textbf{Reverse the boundaries.} Reversing the upper and lower boundaries
proves the single-input assertion.

\item \textbf{Match the whole family.} Vertical maps must be matched for the
whole family; they need not allow the family to be pasted one member at a
time with other members held fixed.\qedhere
\end{enumerate}
\end{proof}

\begin{remark}[Empty profiles]
A cell with two empty boundaries has a specified upper object, a specified
lower object, and two lateral arrows between them. It is an ordinary
possible label of a zero-port rectangle. Its arbitrary vertical composites
are available in a merge calculus; they are not automatically tree
operations. Structural empty squares on a vertical arrow are distinguished
and obey the category laws. These conventions include monad units without
identifying an empty boundary with an atomic horizontal unit.
\end{remark}

\section{Formal categories and their homomorphisms}\label{sec:monads}

\begin{definition}\label{def:formalcat}
A \textbf{formal category} in $\Pcal$ is a quadruple $(X,A,\mu,\eta)$,
where $A:X\harr X$ is an atomic arrow and
$\mu:(A,A)\cell A$, $\eta:\emp X\cell A$ are globular cells such that
\begin{equation}\label{eq:monadlaws}
\begin{tikzcd}[column sep=2.3em]
(A,A,A) \arrow[r,"{(\mu,1_A)}"] \arrow[d,"{(1_A,\mu)}"'] &
(A,A) \arrow[d,"\mu"]\\
(A,A) \arrow[r,"\mu"'] & A
\end{tikzcd}
\qquad
\begin{tikzcd}[column sep=2.2em]
(A,A) \arrow[dr,"\mu"'] &
A \arrow[l,"{(\eta,1_A)}"'] \arrow[r,"{(1_A,\eta)}"]
  \arrow[d,equal] & (A,A) \arrow[dl,"\mu"]\\
& A. &
\end{tikzcd}
\end{equation}
Each expression denotes the corresponding rooted tree evaluation.
A homomorphism $(X,A)\to(Y,B)$ is a vertical arrow $f:X\to Y$ and a
cell $F\in\Pcal_{f,f}(A;B)$ satisfying
\begin{equation}\label{eq:functorlaws}
\begin{tikzcd}
(A,A) \arrow[r,"\mu_A"] \arrow[d,"{(F,F)}"'] & A \arrow[d,"F"]\\
(B,B) \arrow[r,"\mu_B"'] & B
\end{tikzcd}
\qquad
\begin{tikzcd}
\emp X \arrow[r,"\eta_A"] \arrow[d,"{[f]}"'] & A \arrow[d,"F"]\\
\emp Y \arrow[r,"\eta_B"'] & B.
\end{tikzcd}
\end{equation}
Here $[f]$ is the structural empty square on $f$; the notation
$\eta_B[f]$ denotes its paste above $\eta_B$.
\end{definition}

Each complete composite in \cref{eq:monadlaws,eq:functorlaws} is a
tree evaluation. Individual whiskerings, such as $(\mu,1_A)$, are
displayed in the merge realization; adjoining an unused identity wire
is not itself a designated tree operation.

\begin{proposition}\label{prop:cat}
Formal categories and their homomorphisms form a category $\FCat(\Pcal)$.
\end{proposition}
\begin{proof}\leavevmode
\begin{enumerate}
\item \textbf{Define composition.} Compose $(f,F)$ and $(g,G)$ by $(gf,GF)$, using unary tree pasting.
Associativity and identities follow from the corresponding tree equations.

\item \textbf{Preserve multiplication.} For multiplication, the homomorphism equations for $F$ and $G$,
interleaved with substitution associativity, give
\begin{align*}
 (GF)\mu_A
 &=G(F\mu_A)=G\bigl(\mu_B(F,F)\bigr)\\
 &=(G\mu_B)(F,F)=\mu_C(G,G)(F,F)
   =\mu_C(GF,GF).
\end{align*}

\item \textbf{Preserve units.} For units, the structural empty squares compose with their vertical
arrows, so
$(GF)\eta_A=G\eta_B[f]=\eta_C[g][f]=\eta_C[gf]$.
Thus the composite satisfies both homomorphism equations. The same
equations hold for identity cells by the substitution identity laws.\qedhere
\end{enumerate}
\end{proof}

\begin{theorem}[Invariance under the reduct]\label{thm:reduct}
For every merge calculus with designated atomic carriers, the identity on
data induces an isomorphism of categories
\[
 \FCat(\Bcal)\cong\FCat(\Up\Bcal).
\]
\end{theorem}
\begin{proof}\leavevmode
\begin{enumerate}
\item \textbf{Retain the defining cells.} All cells specifying an object or morphism occur in the retained cell
collections.

\item \textbf{Retain the defining equations.} Every operation in \cref{eq:monadlaws,eq:functorlaws} is a
tree or structural operation. The reduct retains their values and the
equalities between them.

\item \textbf{Retain category composition.} Finally, identity and composite homomorphisms use
only structural and unary tree operations. Consequently the object sets,
hom-sets, identities, and compositions agree.\qedhere
\end{enumerate}
\end{proof}

\section{The typed forest envelope}\label{sec:forest}

Let $P$ be a globular virtual calculus: its objects form a set or collection,
its arrows form a directed graph, and $P(\Gamma;H)$ has identities and
substitution for composable paths with the same endpoints. In the main
application $P$ is the entire globular part of $\D$. Arrows between
different objects are retained.

\begin{definition}\label{def:forest}
The strict path $2$-category $\Forest(P)$ has the objects of $P$ and its
finite paths as $1$-arrows. For paths $\Gamma=(H_1,\ldots,H_n)$ and
$\Delta=(K_1,\ldots,K_m)$ with common endpoints, a $2$-cell consists of
a weakly increasing subdivision
\[
 0=i_0\leq i_1\leq\cdots\leq i_m=n
\]
whose cut vertices are the vertices of $\Delta$, and cells
$a_j\in P(H_{i_{j-1}+1},\ldots,H_{i_j};K_j)$.
Repeated indices are permitted only with the indicated coincident vertices.
When $m=0$, there is a cell only if $n=0$, and it is the empty identity
forest. Vertical composition substitutes the $a_j$; horizontal composition
concatenates subdivisions and cells.
\end{definition}

\begin{lemma}\label{lem:forest}
These operations make $\Forest(P)$ a strict $2$-category, and
\begin{equation}\label{eq:forestsingle}
 \Forest(P)(\Gamma;H)=P(\Gamma;H)
\end{equation}
for atomic $H$. A map from $P$ to the single-output part of a strict path
$2$-category extends uniquely to a strict $2$-functor from $\Forest(P)$.
\end{lemma}
\begin{proof}\leavevmode
\begin{enumerate}
\item \textbf{Prove associativity.} Associativity of substitution proves vertical associativity; concatenation
proves horizontal associativity.

\item \textbf{Prove interchange.} Substitution on disjoint consecutive blocks
commutes, which is interchange.

\item \textbf{Verify the identity laws.} The forest of unary identities and the
empty forest give the identity laws.

\item \textbf{Recover single-output cells.} With one output the subdivision has
only one block, proving \cref{eq:forestsingle}.

\item \textbf{Prove the extension property.} A forest must be sent to
the horizontal composite of the images of its roots, so an extension is
unique, and the same substitution identities show that it is a functor.\qedhere
\end{enumerate}
\end{proof}

\begin{definition}[Intrinsic comparison class]\label{def:universal}
A cell $s\in P(\Omega;T)$ is \textbf{contextually universal} if for
every compatible pair of paths $L,R$ and every atomic output $B$, each
$\alpha\in P(L,\Omega,R;B)$ factors uniquely as
\begin{equation}\label{eq:universal}
\begin{tikzcd}[column sep=4em]
(L,\Omega,R) \arrow[r,"1_L\wh s\wh1_R"] \arrow[dr,"\alpha"'] &
(L,T,R) \arrow[d,dashed,"\exists!\,\bar\alpha"]\\
& B.
\end{tikzcd}
\end{equation}
Equivalently, substitution gives a bijection
$s^*:P(L,T,R;B)\to P(L,\Omega,R;B)$.
Let $S(P)$ be the class of all such cells, including those with empty
source. We say that $\Omega$ is \textbf{represented by $T$} when such an
$s$ exists.
\end{definition}

\begin{lemma}\label{lem:universalclosure}
Identities are contextually universal. Substitution of contextually
universal comparisons into a contextually universal comparison is
contextually universal. Contextual instances of such comparisons are
stable under further adjoining contexts.
\end{lemma}
\begin{proof}\leavevmode
\begin{enumerate}
\item \textbf{Include identities.} The first assertion is the identity law.

\item \textbf{Close under substitution.} For the second, the substitution
map of the resulting cell is the composite of the substitution maps for
the outer comparison and its inputs, in their specified contexts. These
maps are bijections.

\item \textbf{Enlarge the context.} For the last assertion, enlarge $L,R$ in
\cref{eq:universal}. This last statement concerns instances in the forest
envelope, which may have several outputs; it does not change the definition
of $S(P)$ as single-output cells.\qedhere
\end{enumerate}
\end{proof}

\begin{lemma}[Globular detection in an equipment]\label{lem:globdet}
If $P$ is the globular part of a virtual equipment $\D$, then
\cref{eq:universal} is equivalent to the same factorization property with
arbitrary exterior vertical boundaries in $\D$.
\end{lemma}
\begin{proof}\leavevmode
\begin{enumerate}
\item \textbf{Restrict the target.} For a target $B$ and exterior arrows $f,g$, its cartesian restriction gives
\[
 \D_{f,g}(\Xi;B)\cong P(\Xi;B[f,g]).
\]

\item \textbf{Transport the factorization test.} This bijection commutes with substitution in $\Xi$. Apply it to both
sides of the desired factorization map and use
\cref{eq:universal} with output $B[f,g]$.

\item \textbf{Recover the globular case.} The converse takes identity
vertical boundaries.\qedhere
\end{enumerate}
\end{proof}

\section{Conservative localization}\label{sec:localization}

\begin{definition}\label{def:localization}
For each $s:\Omega\cell T$ in $S(P)$, adjoin a formal $2$-cell
$s^{-1}:T\cell\Omega$ to $\Forest(P)$. Form the strict $2$-category
presented by these generators, the forest equations, and
\begin{equation}\label{eq:inverserel}
 s^{-1}s=1_\Omega,\qquad ss^{-1}=1_T.
\end{equation}
Denote it by $L(P)$. Equations are closed under horizontal and vertical
composition. Equivalently, localize the hom-categories of $\Forest(P)$
at all contextual instances $1_L\wh s\wh1_R$, and extend horizontal
composition to those localizations.
\end{definition}

The equivalence in the definition follows directly from the universal
property of category localization: horizontal composition sends an
instance to another instance, so it extends uniquely, and its associativity
and interchange follow by uniqueness. In particular, the construction
does not merely invert comparisons in isolated hom-categories and then
postulate compatibility with contexts.

\begin{lemma}[Substitution presheaves]\label{lem:presheaf}
For an atomic arrow $A:X\harr Y$, the assignment
$\Phi_A(\Gamma)=P(\Gamma;A)$ is a presheaf on $\Forest(P)(X,Y)$.
It takes every contextual comparison to a bijection and therefore extends
uniquely to a presheaf $\overline\Phi_A$ on $L(P)(X,Y)$.
The map $P(\Gamma;A)\to L(P)(\Gamma;A)$ is injective.
\end{lemma}
\begin{proof}\leavevmode
\begin{enumerate}
\item \textbf{Define the forest action.} The action of a forest is simultaneous substitution into its consecutive
blocks. Identity and composition follow from the virtual identities.

\item \textbf{Extend across inverses.} The assertion about comparisons is precisely \cref{eq:universal}.
Thus the action of a formal inverse is the inverse bijection, and this
defines the extension.

\item \textbf{Separate original cells.} If original cells $a,b$ have equal images, apply
their presheaf actions to $1_A\in P(A;A)$. The results are $a$ and $b$,
respectively, so $a=b$.\qedhere
\end{enumerate}
\end{proof}

\begin{lemma}[Backward evaluation]\label{lem:backward}
Every $z\in L(P)(\Gamma;A)$ with atomic output $A$ equals the image of
the original cell $\overline\Phi_A(z)(1_A)$.
\end{lemma}
\begin{proof}\leavevmode
\begin{enumerate}
\item \textbf{Choose a layered expression.} Write $z$ as a finite vertical sequence of forests and inverses of
contextual comparisons. Such a sequence exists: distribute horizontal
composites of generators into vertical layers by interchange, adding
identity forests.

\item \textbf{Initialize backward evaluation.} Start at the final boundary $A$ with $1_A$ and traverse
the sequence backwards.

\item \textbf{Pass an ordinary forest.} At an ordinary forest $a:\Xi\cell\Psi$,
replace the current $b\in P(\Psi;A)$ by its substitution $ba$.

\item \textbf{Pass an inverse comparison.} At the inverse of
\[
 t=1_L\wh s\wh1_R:(L,\Omega,R)\cell(L,T,R),
\]
the current cell has source $(L,\Omega,R)$. Let $c$ be its unique factor
through $t$ using \cref{eq:universal}, so that $b=ct$.
Then in the localization $bt^{-1}=c$.

\item \textbf{Complete the induction.} In either case the original cell
now obtained equals the localized suffix that it replaces. Induction over
the finite sequence proves the assertion.

\item \textbf{Prove independence of expression.} These two backward operations
are exactly the actions defining $\overline\Phi_A$, which also proves
independence from the sequence used.\qedhere
\end{enumerate}
\end{proof}

\begin{theorem}[Conservativity]\label{thm:conservative}
The canonical comparison is a bijection
\begin{equation}\label{eq:conservative}
 P(\Gamma;A)\xrightarrow{\ \cong\ }L(P)(\Gamma;A)
\end{equation}
for every original path and atomic output. It respects all original
substitutions. Reversing cells gives the corresponding single-input
statement for the dual construction.
\end{theorem}
\begin{proof}\leavevmode
\begin{enumerate}
\item \textbf{Establish the bijection.} Injectivity is \cref{lem:presheaf} and surjectivity is
\cref{lem:backward}.

\item \textbf{Preserve substitution.} The canonical map respects substitution by its
definition.

\item \textbf{Obtain the dual assertion.} Reversing the displayed bijections proves the dual assertion.\qedhere
\end{enumerate}
\end{proof}

\begin{theorem}[Detection of represented paths]\label{thm:detection}
For an original path $\Gamma$ and atomic arrow $T$,
\[
 \Gamma\cong T\text{ in }L(P)
 \quad\Longleftrightarrow\quad
 \Gamma\text{ is represented by }T\text{ in }P.
\]
More precisely, an original cell $s:\Gamma\cell T$ is invertible in
$L(P)$ if and only if $s\in S(P)$.
\end{theorem}
\begin{proof}\leavevmode
\begin{enumerate}
\item \textbf{Invert universal comparisons.} The reverse implication is the construction.

\item \textbf{Recover an original comparison.} Conversely an isomorphism
$\Gamma\to T$ is the image of a unique original cell $s$ by
\cref{thm:conservative}.

\item \textbf{Verify contextual universality.} Whiskering it by any $L,R$ is an isomorphism.
Precomposition with that isomorphism is a bijection into every atomic
output $B$. Apply \cref{eq:conservative} to both its source profiles;
the resulting bijection is \cref{eq:universal}.

\item \textbf{Test a specified comparison.} The same argument applies
to any specified $s$ whose image is invertible.\qedhere
\end{enumerate}
\end{proof}

\begin{theorem}[Universal property]\label{thm:localUP}
An interpretation of $P$ in the single-output part of a strict path
$2$-category extends to $L(P)$ if and only if it sends every member of
$S(P)$ to an invertible comparison. The extension is unique once its
values on the original signature are fixed.
\end{theorem}
\begin{proof}\leavevmode
\begin{enumerate}
\item \textbf{Extend to forests.} Extend first to forests by \cref{lem:forest}.

\item \textbf{Prove necessity.} Necessity follows from
\cref{eq:inverserel}.

\item \textbf{Define the inverse images.} For sufficiency send $s^{-1}$ to the inverse of
the interpreted $s$; all defining equations hold.

\item \textbf{Prove uniqueness.} An inverse, when it
exists, is unique, so all values of the extension are forced.\qedhere
\end{enumerate}
\end{proof}

\begin{remark}
The backward evaluator is a normalization theorem for cells with one
atomic output. It makes no assertion that arbitrary many-output cells
have a forest normal form. Such an assertion is unnecessary for
\cref{thm:conservative,thm:detection,thm:localUP}.
\end{remark}

\section{Changing vertical boundaries}\label{sec:global}

Fix a virtual equipment $\D$, and let $P$ be its entire globular calculus.
The constructions below use only represented units and restrictions already
present in $\D$. For the displayed calculations, choose representatives for
the finitely many objects and vertical arrows involved and write
$f_*:X\harr X'$ and $f^*:X'\harr X$ for the resulting companion and
conjoint. No simultaneous choice of such representatives is part of the
construction: \cref{prop:choices} below quotients all local presentations by
their unique binding-compatible comparisons. This is the familiar binding or
base-change calculus; compare \cite[Sections 4--5]{Shulman2008}.

\begin{lemma}[Binding and restriction comparisons]\label{lem:binding}
The units and restrictions of $\D$ give the following data in $L(P)$.
\begin{enumerate}
\item For $f:X\to X'$, an adjunction $f_*\dashv f^*$ with
\[
 j_f:\emp X\cell f_*f^*,\qquad
 e_f:f^*f_*\cell\emp{X'},
\]
satisfying
\begin{equation}\label{eq:triangles}
\begin{tikzcd}[column sep=2.2em]
f_* \arrow[r,"j_f\wh1"] \arrow[dr,equal] &
f_*f^*f_* \arrow[d,"1\wh e_f"]\\
& f_*
\end{tikzcd}
\qquad
\begin{tikzcd}[column sep=2.2em]
f^* \arrow[r,"1\wh j_f"] \arrow[dr,equal] &
f^*f_*f^* \arrow[d,"e_f\wh1"]\\
& f^*.
\end{tikzcd}
\end{equation}
\item For $H:X'\harr Y'$ and $f:X\to X'$, $g:Y\to Y'$, an
intrinsically universal comparison, with unit factors suppressed,
\begin{equation}\label{eq:restrictioncomparison}
 s_{f,H,g}:f_*Hg^*\cell H[f,g].
\end{equation}
\item For composable $f:X\to Y$, $u:Y\to Z$, invertible comparisons
\[
 k_{f,u}:f_*u_*\cell(uf)_*,\qquad
 l_{f,u}:u^*f^*\cell(uf)^*.
\]
They are associative and unital, and identify the displayed adjunction
for $uf$ with the composite adjunction.
\end{enumerate}
\end{lemma}
\begin{proof}\leavevmode
\begin{enumerate}
\item \textbf{Construct the binding pairs.} We recall the binding argument, to specify which universal properties are
used. The cartesian cell $f_*\cell U_{X'}$ has sides $(f,1)$.
Factoring the unit cell on $f$ through it gives the other companion
binding cell, with sides $(1,f)$. The corresponding two cells for $f^*$
have the reversed sides. A vertical paste of each pair is the unit square
on $f$.

\item \textbf{Prove the horizontal binding identity.} The horizontal binding identity is proved by postcomposing with
the defining cartesian cell: both sides then give that same unit square,
and uniqueness of cartesian factorization makes them equal.

\item \textbf{Construct the restriction comparison.} Paste the cartesian companion and conjoint cells on the two ends of
$H$, and remove the displayed copies of $U_{X'}$ and $U_{Y'}$ by their
unit universal properties. Factoring through $H[f,g]$ gives
\cref{eq:restrictioncomparison}.

\item \textbf{Prove contextual universality.} To factor a cell having $f_*Hg^*$
in its source, insert the other two binding cells on its sides and then
use the cartesian factorization defining $H[f,g]$. Substitution back gives
the original cell by the horizontal binding identities. In the opposite
order, cartesian uniqueness gives the identity factorization. This works
inside any $L,R$ and with arbitrary exterior maps; hence the comparison
belongs to $S(P)$. This is the binding form of the opcartesian base-change
comparison of \cite[Theorem 7.16]{CruttwellShulman2010}.

\item \textbf{Construct the adjunction unit.} For the first assertion, the same argument with $H=U_{X'}$ makes
$f_*f^*$ a representative of $U_{X'}[f,f]$, after removal of a unit
factor. Let $t_f:\emp X\cell U_{X'}[f,f]$ be the factor of the
nullary unit at $X'$ with sides $(f,f)$. Define $j_f$ by composing
$t_f$ with the inverse representing comparison.

\item \textbf{Construct the adjunction counit.} The two cartesian
binding cells, followed by unit composition, give
$f^*f_*\cell U_{X'}$; compose with the inverse of
$\emp{X'}\cell U_{X'}$ to obtain $e_f$.

\item \textbf{Verify the triangle identities.} The two composites in \cref{eq:triangles}, on expanding these
definitions, are the horizontal binding pastes just proved. Removal and
insertion of the representing comparisons cancel. Thus they are the
identity cells.

\item \textbf{Construct composite binding comparisons.} For composable $f:X\to Y$ and $u:Y\to Z$, paste their companion binding
pairs. Factoring the resulting binding pair through a companion binding
pair for $uf$ gives a comparison
$k_{f,u}:f_*u_*\cell(uf)_*$. By
\cite[Corollary 7.19]{CruttwellShulman2010}, the pasted companion path has
an opcartesian composite comparison to the base-change object of $uf$.
That comparison is characterized by the same pasted binding cells as
$k_{f,u}$, so cartesian uniqueness identifies it with the comparison just
constructed. Its factorization property holds in every surrounding
horizontal context; hence \cref{lem:globdet} places $k_{f,u}$ in $S(P)$,
and it is invertible in $L(P)$. Dually,
the pasted conjoint binding pairs give
$l_{f,u}:u^*f^*\cell(uf)^*$, which likewise belongs to $S(P)$ and is
invertible in $L(P)$.

\item \textbf{Prove comparison coherence.} Two comparisons between binding pairs which respect the binding cells are
equal: postcompose with a cartesian binding cell and use its uniqueness.
For three composable vertical arrows, the two composites of the $k$'s
carry the same pasted companion binding pair to the binding pair of the
threefold composite, so they agree; the same argument applies to the
$l$'s. Identity comparisons are proved similarly.

\item \textbf{Identify composite adjunctions.} Finally the constructions
of $j_f,e_f$ from the binding cells show, again by uniqueness, that these
comparisons identify the displayed adjunction for a composite with the
composite adjunction.\qedhere
\end{enumerate}
\end{proof}

We normalize a local presentation of a vertical identity to the empty path
with identity adjunction. The comparison $\emp X\cong U_X$ permits this
normalization without changing the atomic graph. For later calculations,
compatibility with composition in \cref{lem:binding} says explicitly
\begin{align}
 j_{uf}&=(k_{f,u}\wh l_{f,u})
       (1_{f_*}\wh j_u\wh1_{f^*})j_f,\label{eq:unitcomp}\\
 e_{uf}&=e_u(1_{u^*}\wh e_f\wh1_{u_*})
       (l_{f,u}^{-1}\wh k_{f,u}^{-1}).\label{eq:counitcomp}
\end{align}

\begin{definition}[Cloven global merge presentation]\label{def:global}
Given a system $\mathfrak b$ of represented units, companions, and
conjoints, define an auxiliary presentation $\E_{\mathfrak b}(\D)$.
Its objects and vertical arrows are those of $\D$, and its horizontal
boundaries are paths in the original horizontal graph. Define
\begin{equation}\label{eq:globalcells}
 \E_{\mathfrak b}(\D)_{f,g}(\Gamma;\Delta)
 :=L(P)(\Gamma;f_*\Delta g^*).
\end{equation}
For a square $a$, denote its right-hand representative by $\widehat a$.
Horizontal pasting of $a$ over $(f,g)$ and $b$ over $(g,h)$ is
\begin{equation}\label{eq:hcomp}
\begin{tikzcd}[column sep=4em]
\Gamma\Gamma' \arrow[r,"\widehat a\wh\widehat b"]
  \arrow[dr,"\widehat{a\wh b}"'] &
f_*\Delta g^*g_*\Delta'h^*
  \arrow[d,"1\wh e_g\wh1"]\\
& f_*\Delta\Delta'h^*.
\end{tikzcd}
\end{equation}
If $b$ below $a$ has sides $(u,v)$, define vertical pasting by
\begin{equation}\label{eq:vcomp}
\begin{tikzcd}[column sep=6em]
\Gamma \arrow[r,"\widehat a"] \arrow[d,"\widehat{ba}"'] &
f_*\Delta g^* \arrow[d,"1\wh\widehat b\wh1"]\\
(uf)_*\Xi(vg)^* & f_*u_*\Xi v^*g^*
\arrow[l,"k_{f,u}\wh1_\Xi\wh l_{g,v}"]
\end{tikzcd}
\end{equation}
where $\Xi$ is the lower boundary of $b$.
The vertical identity on $\Gamma$ is $1_\Gamma$; the horizontal identity
on $f$ is $j_f$.
\end{definition}

\begin{theorem}[Pasting coherence]\label{thm:globalcoherence}
For every cloven presentation $\mathfrak b$, the operations of
\cref{def:global} give $\E_{\mathfrak b}(\D)$ the structure of a strict
double category with paths as horizontal arrows, and hence a merge
calculus on the original atomic graph.
\end{theorem}
\begin{proof}\leavevmode
\begin{enumerate}
\item \textbf{Prove horizontal associativity.} For three horizontally adjacent squares, either association tensors their
representatives and contracts the two disjoint strings $g^*g_*$ and
$h^*h_*$. The contractions commute by interchange in $L(P)$.

\item \textbf{Verify horizontal identities.} Horizontal composition with $j_f$ or $j_g$ is an identity by the two
triangle identities \cref{eq:triangles}.

\item \textbf{Prove vertical associativity.} For three vertically adjacent squares, either association first substitutes
all three representatives and then compares the three successive
companions and the three successive conjoints with those of the composite
vertical maps. Associativity of $k,l$ from \cref{lem:binding} makes
these results equal.

\item \textbf{Verify vertical identities.} Their identity comparisons are normalized, proving
the vertical identity laws. Equation~\eqref{eq:unitcomp} states exactly
that the horizontal identity on $uf$ is the vertical paste of those on
$f$ and $u$.

\item \textbf{Prove interchange.} It remains to check interchange. At an interior vertical side, write its
upper and lower segments as $g$ and $v$. After expanding the four square
representatives, either order of evaluation has the same uncontracted
string. Pasting horizontally first contracts $g^*g_*$, then $v^*v_*$,
in that string. Pasting vertically first replaces the string
$v^*g^*g_*v_*$ by $(vg)^*(vg)_*$ using $l,k$, and contracts by
$e_{vg}$. These operations are equal by \cref{eq:counitcomp}.
Every other part of the two expressions agrees by functoriality of
whiskering. This proves interchange. It also proves compatibility with
structural subdivisions and hence evaluation of all rectangular schemes.\qedhere
\end{enumerate}
\end{proof}

\begin{proposition}[Choice-free global merge realization]\label{prop:choices}
For each vertical arrow $f$, let $\mathcal B(f)$ be the collection of all
binding presentations furnished by \cref{lem:binding}, obtained from any
represented unit and the corresponding restrictions. If
$\mathfrak f,\mathfrak f'\in\mathcal B(f)$, there are unique
binding-compatible isomorphisms
\[
 \kappa_{\mathfrak f,\mathfrak f'}:
 \mathfrak f_*\cell\mathfrak f'_*,\qquad
 \lambda_{\mathfrak f,\mathfrak f'}:
 \mathfrak f^*\cell\mathfrak f'^*,
\]
and they satisfy the identity and cocycle equations. The cell collections
\begin{equation}\label{eq:choicefreecells}
 \E(\D)_{f,g}(\Gamma;\Delta)
 :=
 \left(
 \coprod_{\substack{\mathfrak f\in\mathcal B(f)\\
                     \mathfrak g\in\mathcal B(g)}}
 L(P)(\Gamma;\mathfrak f_*\Delta\mathfrak g^*)
 \right)\big/\!\sim
\end{equation}
where
\begin{equation}\label{eq:choicefreerel}
 (\mathfrak f,\mathfrak g,a)\sim
 (\mathfrak f',\mathfrak g',a')
 \quad\Longleftrightarrow\quad
 a'=(\kappa_{\mathfrak f,\mathfrak f'}\wh1_\Delta\wh
        \lambda_{\mathfrak g,\mathfrak g'})a,
\end{equation}
form a merge realization on the original boundary signature. Every finite
rectangular pasting is evaluated by choosing representatives for its
finitely many cells and vertical seams, evaluating by the local cloven
pasting diagrams \eqref{eq:hcomp} and \eqref{eq:vcomp}, and taking the class in
\cref{eq:choicefreecells}. The result is independent of all such local
choices. Thus $\E(\D)$ is defined without a simultaneous choice of units,
companions, or conjoints. Any cloven presentation
$\E_{\mathfrak b}(\D)$ is canonically isomorphic to $\E(\D)$, and the
canonical isomorphisms between cloven presentations satisfy the cocycle
identity.
\end{proposition}
\begin{proof}\leavevmode
\begin{enumerate}
\item \textbf{Compare represented units.} For two represented units, universal factorization gives a unique
isomorphism carrying one nullary universal cell to the other.

\item \textbf{Compare binding pairs.} For two
companion binding pairs, factor the cartesian binding cell of the first
through that of the second and conversely. Each resulting composite has
the same factorization through a cartesian binding cell as the identity,
so cartesian uniqueness makes the two factors inverse. This gives
$\kappa_{\mathfrak f,\mathfrak f'}$; the conjoint argument gives
$\lambda_{\mathfrak f,\mathfrak f'}$.

\item \textbf{Prove the cocycle identities.} A composite of two
binding-compatible comparisons is binding-compatible, so uniqueness gives
\[
 \kappa_{\mathfrak f',\mathfrak f''}
 \kappa_{\mathfrak f,\mathfrak f'}
 =\kappa_{\mathfrak f,\mathfrak f''},\qquad
 \lambda_{\mathfrak f',\mathfrak f''}
 \lambda_{\mathfrak f,\mathfrak f'}
 =\lambda_{\mathfrak f,\mathfrak f''},
\]
together with the identity equations. Hence \cref{eq:choicefreerel} is
an equivalence relation.

\item \textbf{Preserve the binding operations.} Binding-compatible isomorphisms preserve $j$ and $e$: after expanding
these cells by their defining binding factorizations, both possible
transports have the same composite with the relevant cartesian binding
cell, so uniqueness identifies them. They likewise preserve $k$ and $l$,
since the transported and untransported composites respect the same
pasted binding pairs.

\item \textbf{Make pasting independent of choices.} Therefore changing any representative or seam in a
finite rectangular pasting transports its cloven value exactly by the
relation \cref{eq:choicefreerel}. The class of the value is independent
of the choices. Only finitely many representatives are needed for a finite
scheme, so this definition requires no simultaneous cleavage.

\item \textbf{Descend the double-category laws.} Evaluate the two sides of any associativity, identity, or interchange law
using the same finite choices. The calculation in the proof of
\cref{thm:globalcoherence} uses only those finitely many choices and makes
the two values equal; hence they determine the same quotient class. Thus
the quotient is a strict double category and a merge calculus.

\item \textbf{Recover each cloven presentation.} If a
cloven presentation $\mathfrak b$ is given, send
$a\in L(P)(\Gamma;f_*\Delta g^*)$ to its class. Every quotient class has
a unique representative after transport to the $\mathfrak b$-presentation
of its two exterior arrows, so this map is bijective on every cell
collection and respects all pasting. It is therefore a canonical
isomorphism $\E_{\mathfrak b}(\D)\cong\E(\D)$. The cocycle identity is
the one already proved for the binding-compatible comparisons.\qedhere
\end{enumerate}
\end{proof}

\begin{theorem}[Global conservativity]\label{thm:globalembedding}
There is a virtual functor $\D\to\E(\D)$, identical on its boundary
signature, inducing bijections
\begin{equation}\label{eq:globalembedding}
 \D_{f,g}(\Gamma;H)\cong\E(\D)_{f,g}(\Gamma;H)
\end{equation}
for every atomic $H$. It preserves and reflects original cartesian cells.
\end{theorem}
\begin{proof}\leavevmode
\begin{enumerate}
\item \textbf{Select local representatives.} For the fixed profile in question, choose binding presentations of the two
exterior arrows. By \cref{prop:choices}, every class in
$\E(\D)_{f,g}(\Gamma;H)$ has a unique representative in the corresponding
summand of \cref{eq:choicefreecells}.

\item \textbf{Construct the cell bijection.} The desired bijection is therefore
the top edge of
\[
\begin{tikzcd}[column sep=3.2em]
\D_{f,g}(\Gamma;H) \arrow[r,"\cong"] \arrow[d,"\cong"'] &
\E(\D)_{f,g}(\Gamma;H) \arrow[d,"\cong"]\\
P(\Gamma;H[f,g]) \arrow[r,"\cong"'] & L(P)(\Gamma;f_*Hg^*).
\end{tikzcd}
\]
The left edge is cartesian factorization. The bottom edge is the
conservativity bijection of \cref{thm:conservative} followed by
postcomposition with $s_{f,H,g}^{-1}$. The right edge selects the local
representative of the quotient class. The construction
is independent of the two local binding presentations by
\cref{eq:choicefreerel}.

\item \textbf{Preserve identities.} Identities are preserved by the binding identities.

\item \textbf{Preserve virtual substitution.} For substitution of a
compatible family, choose binding presentations for the finitely many
vertical seams and apply the displayed construction to every cell. At each
internal side the two binding cells are adjacent and contract by $e_g$, as
in \cref{eq:hcomp}; at the exterior sides their vertical composites are
compared by $k,l$, as in \cref{eq:vcomp}. The horizontal binding identity
cancels each inserted pair. The remaining cartesian factor is precisely
that of the original substituted cell. Passing to quotient classes proves
compatibility with all virtual substitution, including empty blocks.

\item \textbf{Preserve and reflect cartesianness.} Finally, the test for an original unary cartesian cell is a bijection of
original single-output cell collections for every top path and exterior
boundary. The displayed bijections identify all these tests and their
substitution maps in the two calculi, proving both preservation and
reflection.\qedhere
\end{enumerate}
\end{proof}

\begin{definition}[Equipment boundary transport]\label{def:equipment}
A \textbf{boundary presentation} of a vertical arrow $f:X\to X'$ consists
of arrows $f_*:X\harr X'$ and $f^*:X'\harr X$ (empty paths are permitted
for identities) and cells
\[
 j_f:\emp X\cell f_*f^*,\qquad
 e_f:f^*f_*\cell\emp{X'}
\]
satisfying \cref{eq:triangles}. A poly-virtual calculus has
\textbf{equipment boundary transport} when every vertical arrow admits a
boundary presentation and, for every local choice of presentations for
$f,g$, there are bijections
\[
 \Pcal_{f,g}(\Gamma;\Delta)
   \cong\Pcal_{\mathrm{glob}}(\Gamma;f_*\Delta g^*).
\]
For two local presentations of the same vertical arrow, these
representability bijections induce the unique comparison isomorphisms
between their companion boundaries and between their conjoint boundaries:
with the other exterior map an identity and the lower path empty, the two
companion boundaries represent the same cell functor, and the conjoint
case is dual. The boundary-bending bijections are required to be invariant
under these canonical comparisons, which also carry the corresponding
$j,e$ cells to one another. There are also natural identifications between the cell
collections obtained by iterated boundary bending and by bending along
composite vertical maps. Coherence means that all diagrams of these
bijections generated by inserting identity vertical maps, reassociating a
composable string of vertical maps, changing local boundary presentations
by their canonical comparisons, and the defined tree substitutions
commute. This is coherence of functions on cell collections; it does not
declare a multiwire merge to be a tree operation. A \textbf{poly-virtual
equipment} is a poly-virtual calculus having this property. Transported
boundaries such as $f_*\Delta g^*$ need not have atomic representatives.
In the realizations constructed below the comparison and coherence clauses
are consequences of \cref{lem:binding,prop:choices,prop:transportcoherence},
not additional structure imposed on the construction.
\end{definition}

\begin{proposition}[Mates in the reduct]\label{prop:mates}
For any local choices of binding presentations for $f$ and $g$, the
realization and its reduct have coherent bijections
\begin{equation}\label{eq:mates}
\begin{tikzcd}[column sep=3em]
\E(\D)_{f,g}(\Gamma;\Delta)
  \arrow[r,"\cong"] \arrow[dr,"\cong"'] &
L(P)(\Gamma;f_*\Delta g^*) \arrow[d,"\text{mates}","\cong"']\\
& L(P)(f^*\Gamma g_*;\Delta).
\end{tikzcd}
\end{equation}
The second bijection and its inverse use only single-wire cuts.
Together with \cref{prop:transportcoherence}, these bijections give
$\Up\E(\D)$ equipment boundary transport.
\end{proposition}
\begin{proof}\leavevmode
\begin{enumerate}
\item \textbf{Choose the quotient representative.} The first bijection chooses the unique representative of the quotient class
in the selected summand of \cref{eq:choicefreecells}; equivalently it is
\cref{eq:globalcells} in that cloven presentation.

\item \textbf{Construct the mate.} For the second, cut the
leftmost exposed $f_*$ output against $e_f$, adjoining the input $f^*$;
cut the rightmost $g^*$ output against $e_g$, adjoining $g_*$. In merge
notation the resulting map is
\[
 a\longmapsto
 (e_f\wh1_\Delta\wh e_g)(1_{f^*}\wh a\wh1_{g_*}).
\]

\item \textbf{Construct and verify the inverse.} For its inverse, cut the $f^*$ output of $j_f$ into the leftmost input
and the $g_*$ output of $j_g$ into the rightmost input. Each cut satisfies
\cref{eq:cutplacement}. The two composites reduce to identity by
\cref{eq:triangles}.

\item \textbf{Prove substitution and composition compatibility.} Their naturality under tree substitution follows
by substituting the same cuts into the corresponding trees. Compatibility
with vertical composition is \cref{eq:unitcomp,eq:counitcomp}. All
these operations and equalities are retained by the reduct.\qedhere
\end{enumerate}
\end{proof}

\begin{proposition}[Finite coherence for boundary transport]\label{prop:transportcoherence}
For a composable finite string $f_1,\ldots,f_n$, iterated use of the
comparisons $k$ gives a comparison
\[
 K_{f_1,\ldots,f_n}:
 f_{1*}\cdots f_{n*}\cell(f_n\cdots f_1)_*,
\]
and iterated use of $l$ gives
\[
 L_{f_1,\ldots,f_n}:
 f_n^*\cdots f_1^*\cell(f_n\cdots f_1)^*.
\]
These comparisons are independent of all bracketings. Every iterated
boundary-bending bijection in \cref{def:equipment} reduces to a common
normal form built from $j,e,K,L$, and all the coherence diagrams required
there commute. Hence no further boundary-transport coherence equations
are imposed in the realization.
\end{proposition}
\begin{proof}\leavevmode
\begin{enumerate}
\item \textbf{Normalize composite comparisons.} For three arrows, independence of bracketing is the associativity of $k$
and $l$ proved in \cref{lem:binding}; induction gives the assertion for
any finite string. The identity cases use the normalized empty-path
presentations.

\item \textbf{Expand boundary bending.} Expand an iterated bending into insertions of $j$, cuts
against $e$, and the comparison cells joining successive companion or
conjoint paths. Reassociate the latter to $K$ and $L$. There are four
possible local changes of evaluation order.

\item \textbf{Commute disjoint cuts.} Cuts on disjoint wires commute
by interchange.

\item \textbf{Cancel adjacent unit and counit cells.} Adjacent $j,e$ pairs cancel by the triangle identities
\cref{eq:triangles}.

\item \textbf{Rebracket binding strings.} Rebracketings of companion or conjoint strings agree
by the just-proved independence of $K,L$.

\item \textbf{Combine successive boundary maps.} Replacing bending successively
along $f,u$ by bending along $uf$ is exactly
\cref{eq:unitcomp,eq:counitcomp}.

\item \textbf{Identify all evaluation orders.} Thus any two evaluations with the same
typed exterior reduce to the same expression. These calculations are
spelled out in \cref{app:bindings}.

\item \textbf{Descend through changes of presentation.} Finally, changing the local binding
presentations transports every term by the unique comparisons of
\cref{prop:choices}, so the conclusion descends to the choice-free
realization.\qedhere
\end{enumerate}
\end{proof}

\begin{remark}[Restrictions, extensions, and laxity]\label{rem:lax}
For a path $K:X'\harr Y'$, any local boundary presentations for
$f,g$ give the restriction boundary $f_*Kg^*$. In the corresponding cloven
presentation its cartesian correspondence is \cref{eq:globalcells};
the other bending in \cref{eq:mates} gives extension transport.
For $f:X\to X'$, the operation on endopaths
$R_f(H)=f_*Hf^*$ is lax monoidal, with comparison
\[
 f_*Hf^*f_*Kf^*
 \xrightarrow{\,1\wh e_f\wh1\,}f_*HKf^*
\]
and unit $j_f$. Contractions of disjoint counits prove associativity,
and \cref{eq:triangles} proves unitality. No inverse for this comparison
is supplied or needed. Vertical reversal interchanges the roles of
restriction and extension; the resulting transport remains a formal
boundary even when it has no bicomodule representative.
\end{remark}

\section{The completion theorem and maps of calculi}\label{sec:completion}

\begin{theorem}[Poly-virtual completion]\label{thm:completion}
The construction
$\E_{\mathrm{poly}}(\D)=\Up\E(\D)$ has the properties stated in
\cref{thm:maincompletion}. In particular,
\[
 \FCat(\E_{\mathrm{poly}}(\D))\cong\Mon(\D),
\]
where $\Mon(\D)$ denotes the category of monoids and monoid
homomorphisms in $\D$, and carriers on the left are the original atomic
arrows.
\end{theorem}
\begin{proof}\leavevmode
\begin{enumerate}
\item \textbf{Assemble the calculus and transport.} The tree calculus and transport follow from
\cref{thm:globalcoherence,prop:choices,prop:mates,prop:transportcoherence}.

\item \textbf{Identify the monoid data and equations.} Global conservativity is
\cref{thm:globalembedding}, and identifies all multiplication, unit,
and homomorphism cells. Its compatibility with substitution identifies
all their equations and compositions.

\item \textbf{Pass to the reduct.} Finally apply \cref{thm:reduct}.\qedhere
\end{enumerate}
\end{proof}

\begin{theorem}[Exact extension criterion]\label{thm:functoriality}
Let $F:P\to P'$ be a map of typed globular calculi. It extends, with
the same values on atomic arrows and cells, to a map $L(P)\to L(P')$
if and only if $F(S(P))\subseteq S(P')$.
For a virtual functor between equipments satisfying this criterion,
the extension globalizes canonically to the choice-free merge realizations.
In cloven presentations it is represented by the unique binding-compatible
comparisons of \cref{prop:choices}.
\end{theorem}
\begin{proof}\leavevmode
\begin{enumerate}
\item \textbf{Recognize the extension condition.} By \cref{thm:localUP}, extension is equivalent to invertibility of
every $F(s)$ in $L(P')$. By \cref{thm:detection}, this is precisely
membership in $S(P')$.

\item \textbf{Transport the binding cells.} For globalization, the image of an original unit comparison becomes
invertible. Hence the image binding cells define a companion and conjoint
presentation for the image vertical map in the target realization.

\item \textbf{Compare local presentations.} Compare
this local presentation with any other one by the unique binding-compatible
comparisons of \cref{prop:choices}, and use these at the two ends of the
image of \cref{eq:globalcells}.

\item \textbf{Preserve the two compositions.} Their compatibility with $j,e,k,l$ proves
that \cref{eq:hcomp,eq:vcomp} are respected.

\item \textbf{Prove independence and functoriality.} The cocycle identity makes
the resulting quotient class independent of the local presentations, and
uniqueness proves compatibility with compositions of such functors.\qedhere
\end{enumerate}
\end{proof}

\begin{proposition}[Path presentation]\label{prop:pathpresentation}
A merge calculus on a horizontal graph is equivalently a strict double
category whose horizontal category is the free path category on that
graph, with the graph retained as designated atomic arrows.
\end{proposition}
\begin{proof}\leavevmode
\begin{enumerate}
\item \textbf{Construct a strict double category.} From a merge calculus, regard a boundary path as a horizontal arrow and
use its cell collections as squares. Two-box rectangular evaluations
give the two compositions, and scheme substitution gives associativity,
units, and interchange.

\item \textbf{Evaluate rectangular expressions.} Conversely, compose a rectangular expression
in a strict double category. The double-category laws identify the
generating equivalent schemes, so its value is well defined.

\item \textbf{Check that the constructions are inverse.} These
constructions recover exactly the specified cell collections and pasting
operations in either direction.\qedhere
\end{enumerate}
\end{proof}

This is a presentation equivalence for the strict path realization used
here. It does not identify every poly-virtual calculus with a merge
calculus, nor assert that forgetting merges is invertible.

\section{Formal internal categories in an arbitrary monoidal base}\label{sec:formalinternal}

\begin{definition}\label{def:Q}
The \textbf{bicomodule poly-virtual equipment} is
\begin{equation}\label{eq:Q}
 \Q(\V)=\Up\bigl(\E(\Ccal(\V))^{\vop}\bigr).
\end{equation}
Vertical reversal reverses the vertical category and the upper and
lower boundaries of every square, retaining horizontal orientations.
A \textbf{formal internal category in $\V$} is a formal category in
$\Q(\V)$ on an atomic bicomodule carrier.
\end{definition}

The variance is summarized here; ``primitive'' refers to cells before
localization.
\begin{center}
\small
\begin{tabular}{@{}p{.20\textwidth}p{.35\textwidth}p{.35\textwidth}@{}}
\toprule
Data & $\Ccal(\V)$ & Reversed realization\\
\midrule
Objects & Comonoids & Comonoids\\
Vertical arrows & Opposite comonoid maps & Forward comonoid maps\\
Atomic arrows & Bicomodules & Bicomodules\\
Unary cells & Opposite equivariant maps & Forward equivariant maps\\
Primitive arities & Many inputs, one output & One input, many outputs\\
\bottomrule
\end{tabular}
\end{center}

\begin{theorem}\label{thm:arbitraryV}
For every monoidal category $\V$, \cref{eq:Q} defines a poly-virtual
equipment with the displayed variance. Every primitive single-input
cell, including a unary changing-base cell, is preserved and reflected.
All its cells admit finite expressions in primitive cells, pasting, and
inverses of intrinsically universal comparisons.
\end{theorem}
\begin{proof}\leavevmode
\begin{enumerate}
\item \textbf{Complete and reverse the calculus.} Apply \cref{prop:modinput,lem:bicomodunit,thm:completion}, then
reverse the vertical category and the squares. Under reversal the
original companion and conjoint exchange their roles, as do their unit
and counit.

\item \textbf{Identify equipment transport.} Use the second mate description in \cref{eq:mates} to
rewrite the reversed square correspondence in the form required by
\cref{def:equipment}. The triangle and composition identities are
preserved by this reversal, and \cref{prop:transportcoherence} gives the
required coherent equipment transport.

\item \textbf{Recover primitive cells.} The cell assertion is the
dual of \cref{thm:globalembedding}.

\item \textbf{Verify the finite presentation.} The finite-expression assertion
is \cref{def:localization,def:global}; the latter's adjunction cells
were themselves constructed as finite expressions in
\cref{lem:binding}. The reduct changes none of these collections.\qedhere
\end{enumerate}
\end{proof}

A multiplication is therefore an element of the constructed collection
$\Q(\V)((A,A);A)$, and its associativity is an equality between
specified finite expressions. The carrier remains the given bicomodule
$A$. This proves the existence and preservation assertions of
\cref{thm:main}; \cref{thm:reduct} supplies its assertion about tree
operations. It remains to identify these expressions with actual
cotensor multiplication maps whenever the corresponding boundaries are
represented.

\section{Represented cotensors and Aguiar recovery}\label{sec:recovery}

\begin{definition}[Stable cotensor cone]\label{def:stable}
For $M:C\harr D$ and $N:D\harr E$, a stable cotensor cone is an
equalizer
\begin{equation}\label{eq:cotensorfork}
\begin{tikzcd}[column sep=3.2em,row sep=2.6em]
T \arrow[r,hook,"i_{M,N}"] &
M\otimes N \arrow[r,shift left=.7ex,"\rho_M\otimes1"]
  \arrow[r,shift right=.7ex,"1\otimes\lambda_N"'] & M\otimes D\otimes N\\
& Z \arrow[u,"\phi"'] \arrow[ul,dashed,"\exists!\,\bar\phi"] &
\end{tikzcd}
\end{equation}
which remains an equalizer after applying $K\otimes-\otimes L$ for
every pair of objects $K,L$ of $\V$. Write $T=M\square_D N$.
\end{definition}

This is a property of an existing cone, not a condition on the definition
of $\Q(\V)$. Aguiar calls a monoidal category regular when all
equalizers exist and are preserved in this way
\cite[Definition 2.1.1]{Aguiar1997}.

\begin{lemma}\label{lem:cotensorstructure}
A stable cone~\eqref{eq:cotensorfork} has unique exterior coactions
making $i_{M,N}$ equivariant for the exterior coactions of $M,N$.
It is a $C$--$E$ bicomodule. Its inclusion is a contextually universal
cell $(M,N)\cell T$ of $\Ccal(\V)$.
\end{lemma}
\begin{proof}\leavevmode
\begin{enumerate}
\item \textbf{Construct the induced coactions.} The map $(\lambda_M\otimes1)i_{M,N}$ satisfies the balance equation
in $C\otimes M\otimes N$: this follows by commuting $\lambda_M$
past $\rho_M$ in \cref{eq:bicomodcommute}. Since
$1_C\otimes i_{M,N}$ is the corresponding equalizer, there is a
unique $\lambda_T$ with
\[
\begin{tikzcd}[column sep=3.6em]
T \arrow[r,"i_{M,N}"] \arrow[d,dashed,"\exists!\,\lambda_T"'] &
M\otimes N \arrow[d,"\lambda_M\otimes1"]\\
C\otimes T \arrow[r,hook,"1_C\otimes i_{M,N}"'] & C\otimes M\otimes N.
\end{tikzcd}
\]
Construct $\rho_T$ similarly.

\item \textbf{Verify the bicomodule laws.} Postcomposing the two coassociativity
maps with $1_C\otimes1_C\otimes i_{M,N}$ proves their equality,
because this map is monic; counitality uses $i_{M,N}$ in the same
way. The right coaction and commuting-coaction equations follow by
postcomposing with $i_{M,N}\otimes1_E\otimes1_E$ and
$1_C\otimes i_{M,N}\otimes1_E$, respectively. All these maps are
monic by stability.

\item \textbf{Factor a contextual balanced map.} Now take a balanced map
$b:B\to\bigotimes L\otimes M\otimes N\otimes\bigotimes R$.
Its balance equation at $D$ gives a unique factor
\[
 \bar b:B\longrightarrow\bigotimes L\otimes T\otimes\bigotimes R
\]
through $1\otimes i_{M,N}\otimes1$, by stability.

\item \textbf{Reflect the remaining equations.} For every remaining
balance or exterior equation for $\bar b$, postcompose both sides
with the tensor of $i_{M,N}$ and the additional coaction base. The
defining equations for $\lambda_T,\rho_T$ turn the two resulting
maps into the corresponding equal maps for $b$. The postcomposed map
is monic, so the equation for $\bar b$ holds.

\item \textbf{Prove the universal bijection.} Conversely such a
balanced $\bar b$ gives a balanced $b$ by composition. Uniqueness
is that of the equalizer factorization. This is exactly
\cref{eq:universal} in $\Ccal(\V)$.\qedhere
\end{enumerate}
\end{proof}

Let $\mathbb B(\V)=\E(\Ccal(\V))^{\vop}$ denote the merge
calculus before taking its reduct. Under vertical reversal the cone
becomes a cell
\[
 \iota_{M,N}:T\cell(M,N),
\]
whose underlying map in $\V$ is $i_{M,N}$.

\begin{proposition}[Conversion by one cut]\label{prop:tau}
For a stable cotensor cone the construction supplies an inverse
comparison
\[
 \tau_{M,N}:(M,N)\cell T
\]
in $\mathbb B(\V)$. In its poly-virtual reduct, the operation
\begin{equation}\label{eq:tauuniversal}
 \Q(\V)(L,T,R;\Delta)\longrightarrow
 \Q(\V)(L,M,N,R;\Delta),\qquad
 a\longmapsto a(1_L,\tau_{M,N},1_R)
\end{equation}
is a bijection for all compatible paths. Likewise there is a nullary
comparison $\tau_0:\emp C\cell C$ with this property.
\end{proposition}
\begin{proof}\leavevmode
\begin{enumerate}
\item \textbf{Invert the stable comparison.} By \cref{lem:cotensorstructure}, the comparison in
$\Ccal(\V)$ belongs to $S(P)$. Its image is inverted in $L(P)$;
after reversal its inverse has the stated profile.

\item \textbf{Obtain contextual factorization.} Whiskered
precomposition with it is a bijection in the merge calculus.

\item \textbf{Retain the bijection in the reduct.} The map
in \cref{eq:tauuniversal} is a single cut along $T$, so the reduct
retains this map, its domain and codomain, and their equalities.
Consequently it retains the bijection.

\item \textbf{Include nullary comparisons.} The nullary assertion follows
identically from \cref{lem:bicomodunit}.\qedhere
\end{enumerate}
\end{proof}

The inverse comparison identifies two formal boundaries. The cotensor
inclusion $i_{M,N}$ remains its original map in $\V$. The next example
explains why the input substitution bijection requires the constructed
inverse, even after the original output universal property is known.

\begin{example}[One-sided representation does not suffice; Boolean obstruction]\label{ex:boolean}
This is the two-truth-value instance of the Boolean obstruction used to
distinguish poly-bicategorical cut from merge composition in
\cite[Section 5]{Hadzihasanovic2019}, augmented here by the empty structural
square needed for our calculus. Let horizontal arrows be the two truth
values $\bot<\top$ at one object, and give each profile the empty or
singleton collection
\[
 \mathbb H(\Gamma;\Delta)=
 \begin{cases}
  \{*\},&\Gamma=\Delta=\varnothing,\\
  \{*\},&\bigwedge\Gamma\leq\bigvee\Delta,\\
  \varnothing,&\text{otherwise}.
 \end{cases}
\]
The vertical category is terminal. The empty meet is $\top$ and the
empty join is $\bot$; the first case separately supplies the empty
structural square required by \cref{def:poly}. It cannot be a premise
of a single-wire cut because it has no ports. For the remaining cells,
cuts are valid: if every remaining input is true and every remaining
output is false, the first premise forces the cut formula true and the
second forces it false, a contradiction. Applying this argument to the
wires of a tree defines all tree evaluations. Unary identities exist,
and all evaluation equations hold because each cell collection has at
most one element. We obtain a poly-virtual calculus, with the adjoined
empty structural square affecting none of its nonempty profiles.

The output list $(\top,\bot)$ is contextually represented on the
output side by $\top$, since
$\top\vee\bot=\top$ in every join context. Nevertheless
\[
 \mathbb H(\top,\bot;\bot)=\{*\},\qquad
 \mathbb H(\top;\bot)=\varnothing.
\]
Thus one-sided output representation alone does not give the input
substitution bijection used for multiplication. In our construction
that bijection is \cref{eq:tauuniversal}, obtained from a constructed
inverse and then retained under the reduct. This distinguishes its
proof from an inference based only on a one-sided universal property.
\end{example}

We now compare iterated representatives. Their unique cone comparisons
must supply the associator and unitors used by ordinary cotensor
multiplication; this will identify the formal category equations.

\begin{lemma}[Coherence of represented cotensors]\label{lem:cotensorcoherence}
When the stable cotensors of the pairs being composed exist, all
parenthesizations of a finite bicomodule path represent the same balanced
cone functor. There is a unique comparison between any two of them
commuting with their maps to the unparenthesized tensor. These comparisons
are isomorphisms and satisfy the pentagon and triangle identities.
\end{lemma}
\begin{proof}\leavevmode
\begin{enumerate}
\item \textbf{Identify the triple cone functor.} For a triple, the object $(M\square_D N)\square_E K$, mapped into
$M\otimes N\otimes K$, receives a map from $B$ exactly when a map
from $B$ into that tensor satisfies both adjacent balance equations.
Indeed the first equation uniquely factors it through
$i_{M,N}\otimes1_K$, and the second then becomes the balance
equation for the second cotensor. Monicity after tensoring verifies
this equivalence, including exterior equivariance, as in
\cref{lem:cotensorstructure}.

\item \textbf{Construct the associator.} Reversing the order of these two
factorizations gives $M\square_D(N\square_E K)$ and the same
cone functor. Put $T_L=(M\square_D N)\square_E K$ and
$T_R=M\square_D(N\square_E K)$, with cone maps
$I_L=(i_{M,N}\otimes1)i_{M\square N,K}$ and
$I_R=(1\otimes i_{N,K})i_{M,N\square K}$.
The unique mutually inverse factorizations define the associator by
\[
\begin{tikzcd}[column sep=4em]
T_L \arrow[r,"a_{M,N,K}","\sim"'] \arrow[dr,"I_L"'] &
T_R \arrow[d,"I_R"]\\
& M\otimes N\otimes K.
\end{tikzcd}
\]
The tensor associator is understood on the lower vertex.

\item \textbf{Extend to finite paths.} Induction
establishes the same property for any finite path.

\item \textbf{Prove the pentagon identity.} Two routes around
the pentagon give maps between representatives with the same composite
to the unparenthesized fourfold tensor, so uniqueness makes them equal.

\item \textbf{Identify the stable unit cones.} For units, $\lambda_M:M\to C\otimes M$ is the equalizer of
$\delta_C\otimes1$ and $1\otimes\lambda_M$. If $x:B\to C\otimes M$
equalizes them, apply $\eps_C$ to their first tensor factor to obtain
\[
 x=\lambda_M(\eps_C\otimes1)x.
\]
Counitality proves uniqueness. This formula survives arbitrary tensoring,
so the cone is stable. It gives an isomorphism $C\square_C M\cong M$;
the right coaction gives $M\square_D D\cong M$.

\item \textbf{Prove the triangle identities.} Inserting a unit
factor and removing it by its counit preserves the universal cone.
Thus the two routes in each triangle have the same cone and agree
by uniqueness. This also verifies the nullary cases of the induction.\qedhere
\end{enumerate}
\end{proof}

\begin{theorem}[Fixed-base recovery]\label{thm:aguiar}
For any fixed choices of the stable cotensor representatives used to
present cotensor composition, the formal-category structures on a
$C$--$C$ bicomodule $A$ in $\Q(\V)$ are in bijection with monoid
structures on $A$ in the cotensor monoidal category
$\Bicomod_C(\V)$. Explicitly,
\begin{equation}\label{eq:mutransport}
\begin{tikzcd}[column sep=3em]
(A,A) \arrow[r,"\tau_{A,A}","\sim"'] \arrow[dr,"\mu"'] &
A\square_C A \arrow[d,"m"]\\
& A
\end{tikzcd}
\qquad
\begin{tikzcd}[column sep=2.6em]
\emp C \arrow[r,"\tau_0","\sim"'] \arrow[dr,"\eta"'] &
C \arrow[d,"u"]\\
& A.
\end{tikzcd}
\end{equation}
The inverse formulas are $m=\mu\iota_{A,A}$ and
$u=\eta\tau_0^{-1}$, evaluated in the merge presentation and identified
with actual maps by unary full faithfulness.
In particular this recovers Aguiar internal categories in every regular
monoidal category.
\end{theorem}
\begin{proof}\leavevmode
\begin{enumerate}
\item \textbf{Transport multiplication and unit cells.} By \cref{prop:tau} and unary full faithfulness from
\cref{thm:arbitraryV},
\begin{align*}
 \Q(\V)((A,A);A)&\cong
 \Q(\V)(A\square_C A;A)
 \cong\Bicomod_C(\V)(A\square_C A,A),\\
 \Q(\V)(\emp C;A)&\cong
 \Q(\V)(C;A)
 \cong\Bicomod_C(\V)(C,A).
\end{align*}
These are the formulas~\eqref{eq:mutransport}. Their forward
operations each use one cut. We verify all equations.

\item \textbf{Compare the triple cones.} Write $T=A\square_C A$, $T_L=T\square_C A$, $T_R=A\square_C T$,
and let
\[
 I_L=(\iota_{A,A}\wh1_A)\iota_{T,A}:T_L\cell(A,A,A),\qquad
 I_R=(1_A\wh\iota_{A,A})\iota_{A,T}:T_R\cell(A,A,A).
\]
The defining cone equality of \cref{lem:cotensorcoherence} says
$I_Ra_{A,A,A}=I_L$.

\item \textbf{Identify the multiplication squares.} Functoriality of cotensor maps, defined by
equalizer factorization, says that the squares
\[
\begin{tikzcd}[column sep=3em]
 T_L \arrow[r,"m\square_C1"] \arrow[d,"I_L"'] &
 T \arrow[d,"\iota_{A,A}"]\\
 (A,A,A) \arrow[r,"\mu\wh1"'] &(A,A)
\end{tikzcd}
\qquad
\begin{tikzcd}[column sep=3em]
 T_R \arrow[r,"1\square_Cm"] \arrow[d,"I_R"'] &
 T \arrow[d,"\iota_{A,A}"]\\
 (A,A,A) \arrow[r,"1\wh\mu"'] &(A,A)
\end{tikzcd}
\]
commute in the merge calculus. To see this without presupposing any
monoid laws, compose their upper maps with the cotensor inclusions.
The equalities become the defining factorizations for
$m\square_C1$ and $1\square_Cm$, followed by cancellation of
$\tau_{A,A}\iota_{A,A}=1_T$.

\item \textbf{Preserve and reflect associativity.} Postcomposing with $\mu$ therefore transports the two associativity
trees to
\begin{equation}\label{eq:ordinaryassoc}
 m(m\square_C1_A),\qquad
 m(1_A\square_Cm)a_{A,A,A}:T_L\longrightarrow A.
\end{equation}
Since $I_L$ is invertible in the merge calculus and unary maps embed
faithfully, formal associativity is equivalent to equality of the
actual maps in \cref{eq:ordinaryassoc}.

\item \textbf{Preserve and reflect the unit laws.} Let $\ell_A:C\square_C A\to A$ and $r_A:A\square_C C\to A$
be the unitors from \cref{lem:cotensorcoherence}. Their cone
descriptions identify the left unit tree with
$m(u\square_C1_A)\ell_A^{-1}$ and the right unit tree with
$m(1_A\square_Cu)r_A^{-1}$. Thus the formal unit laws are precisely
\[
 m(u\square_C1_A)=\ell_A,\qquad
 m(1_A\square_Cu)=r_A.
\]
This proves preservation and reflection of the full monoid structure.

\item \textbf{Identify internal categories.} Aguiar's internal category is a monoid in this bicomodule monoidal
category \cite[Theorem 2.2.1 and Definition 2.3.1]{Aguiar1997};
regularity supplies every stable cone used above.\qedhere
\end{enumerate}
\end{proof}

The recovery proof is local to its displayed stable paths. If all finite
powers of $A$ are stably represented, their chosen representatives,
including $C$ for the empty power, form a cotensor monoidal category with
all bicomodule maps between them. The cone factorizations of
\cref{lem:cotensorcoherence} identify the cotensor of the $m$th and $n$th
representatives with the $(m+n)$th and supply its associator and unitors.
The same proof therefore gives the correspondence in this generated
category. The changing-base and action comparisons below likewise use
only the stable paths displayed in their proofs.

\begin{corollary}[Represented paths globally]\label{cor:globalrepresented}
When every bicomodule path is stably represented, composition of its
boundaries identifies every completed square with an ordinary
equivariant bicomodule map between the represented boundaries.
This identification respects mixed bases and all pasting.
\end{corollary}
\begin{proof}\leavevmode
\begin{enumerate}
\item \textbf{Represent the square by an actual map.} Replace each boundary by its representative using the invertible cone
comparisons. This gives a unary square, and unary full faithfulness
identifies it with an actual equivariant map.

\item \textbf{Prove independence of replacements.} Independence from
parenthesization and from the sequence of replacements follows from
\cref{lem:cotensorcoherence}.

\item \textbf{Identify composition over the same boundaries.} For adjacent squares, pasting the
representing cones gives the representing cone of their composite;
uniqueness therefore makes composition agree. The exterior vertical
maps have not been changed during any replacement.\qedhere
\end{enumerate}
\end{proof}

\section{Changing-base functors and bimodule actions}\label{sec:changing}

\begin{lemma}[Cotensor maps over a base map]\label{lem:changebase}
Let $F:M\to M'$ and $G:N\to N'$ be equivariant over comonoid maps
$(f,g)$ and $(g,h)$, respectively. For stable cotensors there is a
unique equivariant map
\[
 F\square_g G:M\square_D N\longrightarrow M'\square_{D'}N'
\]
with
\begin{equation}\label{eq:changefactor}
\begin{tikzcd}[column sep=3.5em]
M\square_D N \arrow[r,hook,"i_{M,N}"]
  \arrow[d,dashed,"\exists!\,F\square_g G"'] &
M\otimes N \arrow[d,"F\otimes G"]\\
M'\square_{D'}N' \arrow[r,hook,"i_{M',N'}"'] & M'\otimes N'.
\end{tikzcd}
\end{equation}
These maps respect identities and composition of all three base maps.
\end{lemma}
\begin{proof}\leavevmode
\begin{enumerate}
\item \textbf{Construct the balance factorization.} Using equivariance, the two balance maps of
$(F\otimes G)i_{M,N}$ in \cref{eq:changefactor} are
\[
 (F\otimes g\otimes G)(\rho_M\otimes1)i_{M,N},\qquad
 (F\otimes g\otimes G)(1\otimes\lambda_N)i_{M,N}.
\]
They agree. The equalizer property gives the unique factorization.

\item \textbf{Verify exterior equivariance.} Postcompose its proposed exterior equivariance equations with
$1\otimes i_{M',N'}$ or $i_{M',N'}\otimes1$. The results follow
from exterior equivariance of $F,G$; stability makes these testing maps
monic.

\item \textbf{Preserve identity maps.} Identity maps satisfy the characterization, so their factor is
the identity.

\item \textbf{Preserve composition.} For successive pairs $(F,G)$ and $(F',G')$, either
factorization has composite with the final inclusion
$(F'F\otimes G'G)i_{M,N}$, so uniqueness proves
\[
 (F'\square_{g'}G')(F\square_g G)
       =(F'F)\square_{g'g}(G'G).\qedhere
\]
\end{enumerate}
\end{proof}

\begin{theorem}[Global category recovery]\label{thm:globalaguiar}
In the represented cotensor setting of \cref{thm:aguiar}, the
correspondence is an isomorphism of categories, including homomorphisms
over arbitrary comonoid maps. Such a homomorphism is exactly a comonoid
map $f:C\to D$ and an equivariant map $F:A\to B$ satisfying
\begin{equation}\label{eq:actualfunctors}
 Fu_A=u_Bf,\qquad
 Fm_A=m_B(F\square_f F).
\end{equation}
\end{theorem}
\begin{proof}\leavevmode
\begin{enumerate}
\item \textbf{Identify the carrier map.} Unary global full faithfulness from \cref{thm:arbitraryV} identifies
the formal carrier square with exactly the equivariant map $F$.

\item \textbf{Transport the multiplication equation.} For multiplication, replace the upper boundary $(A,A)$ by
$A\square_C A$. Naturality of the representing comparisons is
\cref{eq:changefactor}. Consequently the two formal homomorphism
trees become the two maps in
\[
\begin{tikzcd}[column sep=3em]
 A\square_C A \arrow[r,"m_A"] \arrow[d,"F\square_f F"'] &
 A \arrow[d,"F"]\\
 B\square_D B \arrow[r,"m_B"'] &B.
\end{tikzcd}
\]

\item \textbf{Transport the unit equation.} For units, the primitive map $f:C\to D$ intertwines the regular
coactions and satisfies $\eps_Df=\eps_C$. Thus it is the map
between the nullary representatives over $f$, and the unit tree
becomes the left equality in \cref{eq:actualfunctors}.

\item \textbf{Identify the functor category.} Full faithfulness reflects both equations. The identity and composite
comparisons follow from \cref{lem:changebase}, so this is an
isomorphism of categories. These are the internal-functor equations
in the cotensor formulation \cite[Section 4.1]{Aguiar1997}.\qedhere
\end{enumerate}
\end{proof}

\begin{definition}[Atomic bimodule action data]\label{def:bimodule}
Let $(X,A)$ and $(Y,B)$ be formal categories. An $A$--$B$ bimodule
is an atomic arrow $M:X\harr Y$ with action cells
\[
 p:(A,M)\cell M,\qquad q:(M,B)\cell M
\]
satisfying the two associativity squares
\begin{equation}\label{eq:leftaction}
\begin{tikzcd}[column sep=2.4em]
(A,A,M) \arrow[r,"{(\mu_A,1)}"] \arrow[d,"{(1,p)}"'] &
(A,M) \arrow[d,"p"]\\
(A,M) \arrow[r,"p"'] & M
\end{tikzcd}
\qquad
\begin{tikzcd}[column sep=2.4em]
(M,B,B) \arrow[r,"{(q,1)}"] \arrow[d,"{(1,\mu_B)}"'] &
(M,B) \arrow[d,"q"]\\
(M,B) \arrow[r,"q"'] & M,
\end{tikzcd}
\end{equation}
the unit triangles
\begin{equation}\label{eq:rightaction}
\begin{tikzcd}[column sep=3em]
(A,M) \arrow[dr,"p"'] &
M \arrow[l,"{(\eta_A,1)}"'] \arrow[r,"{(1,\eta_B)}"]
  \arrow[d,equal] &
(M,B) \arrow[dl,"q"]\\
& M &
\end{tikzcd}
\end{equation}
and the commuting-action square
\begin{equation}\label{eq:commutingactions}
\begin{tikzcd}
(A,M,B) \arrow[r,"{(p,1)}"] \arrow[d,"{(1,q)}"'] &
(M,B) \arrow[d,"q"]\\
(A,M) \arrow[r,"p"'] & M.
\end{tikzcd}
\end{equation}
An equivariant bimodule map over homomorphisms $(f,F):(X,A)\to(X',A')$
and $(g,G):(Y,B)\to(Y',B')$ is a square
$\theta\in\Pcal_{f,g}(M;M')$ with
\begin{equation}\label{eq:modulemap}
\begin{tikzcd}
(A,M) \arrow[r,"p"] \arrow[d,"{(F,\theta)}"'] & M \arrow[d,"\theta"]\\
(A',M') \arrow[r,"p'"'] & M'
\end{tikzcd}
\qquad
\begin{tikzcd}
(M,B) \arrow[r,"q"] \arrow[d,"{(\theta,G)}"'] & M \arrow[d,"\theta"]\\
(M',B') \arrow[r,"q'"'] & M'.
\end{tikzcd}
\end{equation}
\end{definition}

\begin{theorem}[Bimodule comparison]\label{thm:bimodules}
The reduct preserves and reflects bimodules, their equivariant maps, and
composition of these maps. In the represented bicomodule setting they
are exactly bicomodules with cotensor action maps
\[
 p_0:A\square_C M\to M,\qquad
 q_0:M\square_D B\to M
\]
satisfying the associative, unital, and commuting action equations.
Equivariant maps satisfy
\[
 \theta p_0=p'_0(F\square_f\theta),\qquad
 \theta q_0=q'_0(\theta\square_g G).
\]
\end{theorem}
\begin{proof}\leavevmode
\begin{enumerate}
\item \textbf{Retain the formal bimodule equations.} Every expression in
\cref{eq:leftaction,eq:rightaction,eq:commutingactions,eq:modulemap}
is a tree evaluation, so the first assertion follows exactly as in
\cref{thm:reduct}. Composition of maps respects these equations by
substitution associativity.

\item \textbf{Represent the two actions.} For the second assertion, apply \cref{prop:tau} to $A,M$ and $M,B$
and then use unary full faithfulness.

\item \textbf{Transport action associativity.} The left associativity equation,
tested on $(A\square_C A)\square_C M$, becomes
\[
 p_0(m_A\square_C1_M)
   =p_0(1_A\square_Cp_0)a_{A,A,M}.
\]
The right associativity equation is the same comparison for $M,B,B$.

\item \textbf{Transport commuting actions.} The commuting equation, tested on $(A\square_C M)\square_D B$,
becomes
\[
 q_0(p_0\square_D1_B)
   =p_0(1_A\square_Cq_0)a_{A,M,B}.
\]

\item \textbf{Transport action units.} The unit equations become the left and right cotensor unitors, by the
nullary calculation in \cref{thm:aguiar}.

\item \textbf{Transport bimodule maps.} Finally
\cref{lem:changebase} transports \cref{eq:modulemap} to the two
displayed equivariant-map equations. Every step is reversible.\qedhere
\end{enumerate}
\end{proof}

This proves the recovery assertions of \cref{thm:main}, including
changing-base functors and bimodule actions. These actions supply the
atomic distributor carriers in the construction that follows. To obtain
restrictions along arbitrary formal functors, we also use the finite
bicomodule paths already present in the merge realization. The next three
sections construct their balanced multicells, their many-sided completion,
and the representation of distributor composition.

\section{The virtual equipment of formal distributors}\label{sec:distributors}

The action data of \cref{def:bimodule} already describe distributors on
actual bicomodules. To put formal categories at the objects of an
equipment, we must also construct distributor multicells and restrictions
along formal functors. The merge realization supplies the carriers needed
for those restrictions. Throughout this section,
\[
 \mathbb B(\V)=\E(\Ccal(\V))^{\vop},\qquad
 \Q(\V)=\Up\mathbb B(\V).
\]
We use all horizontal paths of $\mathbb B(\V)$ as carriers while keeping
the objects of the new equipment exactly the formal categories already
defined on atomic bicomodules.

\subsection{The carrier equipment}

\begin{definition}[Paths as carriers]\label{def:carrierequipment}
Let $\Kcal(\V)$ have the objects and vertical arrows of $\mathbb B(\V)$,
and let each finite bicomodule path be an individual horizontal arrow of
$\Kcal(\V)$. Set
\begin{equation}\label{eq:carriercells}
 \Kcal(\V)_{f,g}(H_1,\ldots,H_n;K)
 =\mathbb B(\V)_{f,g}(H_1\cdots H_n;K).
\end{equation}
Substitution on the left is the corresponding rectangular evaluation on
the right. The underlying path of an empty list of carriers is the empty
bicomodule path at its specified vertex.
\end{definition}

A cell of \cref{eq:carriercells} has one output \textbf{carrier}, even
when that carrier contains several bicomodules. This use of the path
presentation is part of the construction. For example, composition along
a carrier with several bicomodules is a merge operation in the original
graph and is an ordinary substitution in $\Kcal(\V)$.

\begin{lemma}\label{lem:carrierequipment}
The calculus $\Kcal(\V)$ is a virtual equipment. Its horizontal unit
at $X$ is the carrier $\emp X$. For any local boundary presentations of
$f:X\to X'$ and $g:Y\to Y'$, a restriction of a carrier
$K:X'\harr Y'$ has carrier
\begin{equation}\label{eq:carrierrestriction}
 K[f,g]=f_*Kg^*.
\end{equation}
\end{lemma}
\begin{proof}\leavevmode
\begin{enumerate}
\item \textbf{Establish substitution and represented units.} The strict double-category laws of \cref{thm:globalcoherence}, retained
under vertical reversal, give the virtual substitution laws in
\cref{eq:carriercells}. Concatenation represents every list of carriers,
including the empty list, so $\emp X$ is the represented unit.

\item \textbf{Construct the restriction cell.} Use the reversed companion and conjoint data from
\cref{thm:arbitraryV}. Boundary bending gives
\[
 \mathbb B_{f,g}(H;K)
 \cong\mathbb B_{\mathrm{glob}}(H;f_*Kg^*).
\]
The cell corresponding to $1_{f_*Kg^*}$ is the restriction cell
$\chi_{f,K,g}:f_*Kg^*\cell K$, displayed as
\[
\begin{tikzcd}[column sep=5em,row sep=3.4em]
X \arrow[r,proarrow,"f_*Kg^*",""{name=top,below}] \arrow[d,"f"'] &
Y \arrow[d,"g"]\\
X' \arrow[r,proarrow,"K"',""{name=bottom,above}] & Y'
\arrow[Rightarrow,from=top,to=bottom,shorten <=2pt,shorten >=2pt,"\chi_{f,K,g}"]
\end{tikzcd}
\]

\item \textbf{Verify the full cartesian property.} To verify its full cartesian property,
allow exterior arrows $r:Z\to X$ and $s:W\to Y$. The comparison
maps for composite companions and conjoints identify
\begin{align*}
 \mathbb B_{r,s}(H;f_*Kg^*)
 &\cong\mathbb B_{\mathrm{glob}}(H;r_*f_*Kg^*s^*)\\
 &\cong\mathbb B_{\mathrm{glob}}(H;(fr)_*K(gs)^*)\\
 &\cong\mathbb B_{fr,gs}(H;K).
\end{align*}
Under these bijections the map is postcomposition with $\chi_{f,K,g}$.
The same calculation applies when $H$ is a concatenation of any list of
carriers. This is cartesianness in $\Kcal(\V)$.\qedhere
\end{enumerate}
\end{proof}

\subsection{Distributors and balanced multicells}

\begin{definition}[Formal distributor]\label{def:formaldistributor}
Let $\mathsf A=(X,A,\mu_A,\eta_A)$ and
$\mathsf B=(Y,B,\mu_B,\eta_B)$ be formal categories. A formal
distributor $M:\mathsf A\harr\mathsf B$ has a carrier
$H_M:X\harr Y$ in $\Kcal(\V)$ and cells
\[
 p_M:(A,H_M)\cell H_M,\qquad
 q_M:(H_M,B)\cell H_M
\]
satisfying the associativity, unit, and commuting-action diagrams
\eqref{eq:leftaction}--\eqref{eq:commutingactions}, with $M$ replaced
by $H_M$ and substitution evaluated in $\Kcal(\V)$.
We call it \textbf{bicomodule-carried} when $H_M$ consists of one actual
bicomodule.
\end{definition}

For a bicomodule-carried distributor these complete diagrams are tree
evaluations, so the data and laws are exactly those of \cref{def:bimodule}.
For general path carriers they are merge evaluations in $\mathbb B(\V)$.

\begin{definition}[Balanced distributor multicell]\label{def:distcell}
Consider distributors
$M_i:\mathsf A_{i-1}\harr\mathsf A_i$ for $1\leq i\leq n$,
$N:\mathsf B_0\harr\mathsf B_1$, and formal functors
\[
 \mathsf F=(f,F):\mathsf A_0\to\mathsf B_0,
 \qquad
 \mathsf G=(g,G):\mathsf A_n\to\mathsf B_1.
\]
For $n>0$, put $\Gamma=(H_{M_1},\ldots,H_{M_n})$. A distributor
multicell over $(\mathsf F,\mathsf G)$ is a cell
$\alpha\in\Kcal(\V)_{f,g}(\Gamma;H_N)$ satisfying the exterior squares, where $A_i,B_j$ denote category carriers:
\begin{equation}\label{eq:distexterior}
\begin{tikzcd}[column sep=2.4em]
(A_0,\Gamma) \arrow[r,"{(p_{M_1},1,\ldots,1)}"]
  \arrow[d,"{(F,\alpha)}"'] & \Gamma \arrow[d,"\alpha"]\\
(B_0,H_N) \arrow[r,"p_N"'] & H_N
\end{tikzcd}
\qquad
\begin{tikzcd}[column sep=2.4em]
(\Gamma,A_n) \arrow[r,"{(1,\ldots,1,q_{M_n})}"]
  \arrow[d,"{(\alpha,G)}"'] & \Gamma \arrow[d,"\alpha"]\\
(H_N,B_1) \arrow[r,"q_N"'] & H_N.
\end{tikzcd}
\end{equation}
For an internal interface, let $\Gamma_i^+$ be $\Gamma$ with the
category carrier $A_i$ inserted between $H_{M_i}$ and $H_{M_{i+1}}$.
Let $r_i$ apply $q_{M_i}$ and $\ell_i$ apply $p_{M_{i+1}}$ at that
interface. Interior balance is
\begin{equation}\label{eq:distbalance}
\begin{tikzcd}
\Gamma_i^+ \arrow[r,"r_i"] \arrow[d,"\ell_i"'] &
\Gamma \arrow[d,"\alpha"]\\
\Gamma \arrow[r,"\alpha"'] & H_N.
\end{tikzcd}
\qquad(1\leq i<n).
\end{equation}
For $n=0$, the functors have common domain $\mathsf A=(X,A)$ and
the underlying cell is $\alpha:\emp X\cell H_N$ over $(f,g)$.
The nullary balance square is
\begin{equation}\label{eq:distnullary}
\begin{tikzcd}
A \arrow[r,"{(F,\alpha)}"] \arrow[d,"{(\alpha,G)}"'] &
(B_0,H_N) \arrow[d,"p_N"]\\
(H_N,B_1) \arrow[r,"q_N"'] & H_N.
\end{tikzcd}
\end{equation}
\end{definition}

Let $\Ddist(\V)$ denote this data, with formal categories as objects,
formal functors as vertical arrows, and formal distributors as horizontal
arrows. Its multicell collections are denoted
$\Ddist(\V)_{\mathsf F,\mathsf G}(M_1,\ldots,M_n;N)$.

\begin{proposition}[Closure under substitution]\label{prop:distvirtual}
The specified data form a virtual double category, and there is an
identification
\begin{equation}\label{eq:distMod}
 \Ddist(\V)
 =\left.\Mod(\Kcal(\V))\right|_{\FCat(\Q(\V))}.
\end{equation}
The restriction in \cref{eq:distMod} is full on the displayed objects:
all distributors between them are retained. Its vertical category is
exactly $\FCat(\Q(\V))$.
\end{proposition}
\begin{proof}\leavevmode
\begin{enumerate}
\item \textbf{Check balance within a substituted block.} Define substitution by that of $\Kcal(\V)$. Consider an inserted
category action in a substituted cell. If its position is internal to
one substituted block, the balance equation for that block identifies the
two results.

\item \textbf{Check nonempty interfaces.} At the interface between two nonempty blocks, their exterior
equations move the action to the corresponding inputs of the outer cell;
the outer balance equation then identifies the results.

\item \textbf{Check empty blocks.} At an empty block,
\cref{eq:distnullary} identifies the actions through its two boundary
functors. Repeating this step moves the action past any finite string of
empty blocks. It also covers the case in which every block is empty.

\item \textbf{Check exterior equivariance.} At the two exterior boundaries, the same calculation uses the exterior
equations of the outer cell and composition of the boundary functors.
Thus the substituted cell satisfies all the defining equations.

\item \textbf{Verify the virtual-category laws.} Unary identity cells satisfy exterior equivariance, and have no interior
balance equations. Associativity and identities now follow from the
corresponding identities in $\Kcal(\V)$.

\item \textbf{Identify the module construction.} The definition of a monoid in
$\Kcal(\V)$ on an original atomic bicomodule is exactly
\cref{def:formalcat}: the comparison with $\Q(\V)$ is
\cref{thm:reduct}. Its homomorphisms are exactly the same pairs $(f,F)$
and equations. Bimodules and their multicells in the module construction
are the displayed action, exterior, and balance equations, proving
\cref{eq:distMod}. This is the explicit instance of Leinster's monoids-and-modules
construction \cite[Section 5.3]{Leinster2004}, in the formulation of
\cite[Definition 2.8]{CruttwellShulman2010}, used here.\qedhere
\end{enumerate}
\end{proof}

\subsection{Regular distributors and restrictions}

\begin{lemma}[The regular distributor]\label{lem:distunit}
For a formal category $\mathsf A$, the distributor
\[
 U_{\mathsf A}=(A,\mu_A,\mu_A)
\]
represents the empty distributor path at $\mathsf A$ in $\Ddist(\V)$.
Its universal cell is the original unit $\eta_A$.
\end{lemma}
\begin{proof}\leavevmode
\begin{enumerate}
\item \textbf{Verify the regular distributor.} The monad laws make the displayed data a distributor, and the two unit
laws give \cref{eq:distnullary} for $\eta_A$.

\item \textbf{Insert a unit beside a nonempty context.} We prove the contextual
universal property directly. Let $\alpha:(L,R)\cell N$ be a
distributor multicell whose two paths meet at $\mathsf A$. If $L$ is
nonempty, insert $U_{\mathsf A}$ by the right action on its last
distributor. If $R$ is nonempty, use instead the left action on its first
distributor. Interior balance for $\alpha$ says that the two formulas
agree when both are available.

\item \textbf{Check the inserted equations.} Associativity of the adjacent action
gives balance against the regular action on $U_{\mathsf A}$, and the
remaining balance and exterior equations are inherited from $\alpha$.

\item \textbf{Treat empty contexts and endpoints.} If both contexts are empty, write the boundary functors as
$\mathsf F=(f,F)$ and $\mathsf G=(g,G)$. Define the inserted cell by
\[
 \beta=p_N(F,\alpha)=q_N(\alpha,G).
\]
The equality is \cref{eq:distnullary}. The functor multiplication laws
and the two action associativities prove the exterior equations for
$\beta$. Its two possible expressions prove compatibility with both
exterior actions. When exactly one context is empty, the corresponding
exterior equation supplies the same endpoint calculation.

\item \textbf{Recover the original multicell.} Substitution of $\eta_A$ into the inserted factor recovers $\alpha$
by the action unit laws.

\item \textbf{Prove uniqueness in every context.} Conversely, for a balanced cell $\beta$ with
an inserted $U_{\mathsf A}$, its balance equation at an adjacent
distributor, followed by substitution of $\eta_A$, forces the insertion
formula just given. If both contexts are empty, exterior equivariance
and $\mu_A(1,\eta_A)=1_A$ give
$\beta=p_N(F,\beta\eta_A)$; the right formula follows similarly.
This proves uniqueness in every context and over every exterior pair
of functors.\qedhere
\end{enumerate}
\end{proof}

\begin{theorem}[Distributor restriction]\label{thm:distrestriction}
Let $N:\mathsf A'\harr\mathsf B'$ be a distributor, and let
\[
 \mathsf F=(f,F):\mathsf A\to\mathsf A',\qquad
 \mathsf G=(g,G):\mathsf B\to\mathsf B'
\]
be formal functors. Its restriction $N[\mathsf F,\mathsf G]$ exists
in $\Ddist(\V)$. For any local boundary presentations of $f,g$, it may
be chosen with carrier
\begin{equation}\label{eq:distrestrictioncarrier}
 H=f_*H_Ng^*.
\end{equation}
For the ambient cartesian cell $\chi:H\cell H_N$ over $(f,g)$,
its actions are uniquely specified by
\begin{equation}\label{eq:distrestrictionactions}
\begin{tikzcd}
(A,H) \arrow[r,dashed,"\exists!\,p_H"] \arrow[d,"{(F,\chi)}"'] &
H \arrow[d,"\chi"]\\
(A',H_N) \arrow[r,"p_N"'] & H_N
\end{tikzcd}
\qquad
\begin{tikzcd}
(H,B) \arrow[r,dashed,"\exists!\,q_H"] \arrow[d,"{(\chi,G)}"'] &
H \arrow[d,"\chi"]\\
(H_N,B') \arrow[r,"q_N"'] & H_N.
\end{tikzcd}
\end{equation}
These restrictions are coherent under identities and iterated restriction.
\end{theorem}
\begin{proof}\leavevmode
\begin{enumerate}
\item \textbf{Construct the restricted actions.} \Cref{lem:carrierequipment} supplies $H$ and $\chi$.
Cartesian factorization of the two right sides of
\cref{eq:distrestrictionactions} constructs unique globular actions
on $H$.

\item \textbf{Prove left associativity.} To prove their laws it suffices, by the same cartesian
uniqueness, to postcompose with $\chi$. Left associativity is the
calculation
\begin{align*}
 \chi p_H(\mu_A,1)
 &=p_N(F\mu_A,\chi)\\
 &=p_N\bigl(\mu_{A'}(F,F),\chi\bigr)\\
 &=p_N\bigl(F,p_N(F,\chi)\bigr)
 =\chi p_H(1,p_H).
\end{align*}

\item \textbf{Verify the left unit law.} For the left unit, the functor unit equation and the action unit equation
give
\[
 \chi p_H(\eta_A,1)
 =p_N(F\eta_A,\chi)
 =p_N(\eta_{A'}[f],\chi)=\chi.
\]

\item \textbf{Verify the right action laws.} Right associativity and its unit equation use $G$ in the same way.

\item \textbf{Prove that the actions commute.} For commuting actions,
\begin{align*}
 \chi q_H(p_H,1)
 &=q_N\bigl(p_N(F,\chi),G\bigr)\\
 &=p_N\bigl(F,q_N(\chi,G)\bigr)
 =\chi p_H(1,q_H).
\end{align*}
Thus $H$ is a distributor, and \cref{eq:distrestrictionactions} says
that $\chi$ is its equivariant distributor map over
$(\mathsf F,\mathsf G)$.

\item \textbf{Verify distributor cartesianness.} To verify cartesianness in $\Ddist(\V)$, take a distributor multicell
$\alpha$ with target $N$ whose exterior functors factor through
$\mathsf F,\mathsf G$. Its unique underlying factor through $\chi$
exists in $\Kcal(\V)$. Postcomposing each proposed exterior equation
for the factor with $\chi$ gives the corresponding equation for
$\alpha$, using \cref{eq:distrestrictionactions}. Postcomposing its
interior balance equations gives those for $\alpha$. In the nullary
case the same argument gives \cref{eq:distnullary}. Cartesian uniqueness
reflects these equalities, so the factor is a distributor multicell.
Its uniqueness is already the underlying cartesian uniqueness.

\item \textbf{Compare successive restrictions.} For further functors $\mathsf R=(r,R)$ and $\mathsf S=(s,S)$,
the composite ambient cartesian cell identifies the carriers
$r_*f_*H_Ng^*s^*$ and $(fr)_*H_N(gs)^*$ by the companion and conjoint
comparisons. Both induced distributor actions are uniquely characterized
by their composites to $N$, with boundary functors
$\mathsf F\mathsf R$ and $\mathsf G\mathsf S$. Therefore the
comparison is an invertible distributor map.

\item \textbf{Prove restriction coherence.} The identity comparisons and
every association comparison respect the same cartesian cells; uniqueness
proves their coherence.\qedhere
\end{enumerate}
\end{proof}

\begin{corollary}\label{cor:distequipment}
For every monoidal $\V$, $\Ddist(\V)$ is a virtual equipment with
vertical category $\FCat(\Q(\V))$.
\end{corollary}
\begin{proof}\leavevmode
\begin{enumerate}
\item \textbf{Assemble the virtual equipment.} Combine \cref{prop:distvirtual,lem:distunit,thm:distrestriction}.
These constructions also verify explicitly the units and restrictions
provided abstractly for $\Mod$ by
\cite[Propositions 5.5 and 7.4]{CruttwellShulman2010}.\qedhere
\end{enumerate}
\end{proof}

The objects selected in \cref{eq:distMod} have not been enlarged:
their carriers are still actual bicomodules. The horizontal carriers
include the paths \cref{eq:distrestrictioncarrier}, with their actions
constructed above. Closure under these restrictions is established inside
the already defined merge realization.

\section{The poly-virtual equipment of formal distributors}\label{sec:distributorcompletion}

\begin{definition}\label{def:FDist}
The poly-virtual equipment of formal distributors is
\begin{equation}\label{eq:FDist}
 \FDist(\V)=\Up\E(\Ddist(\V)).
\end{equation}
Its designated atomic horizontal arrows are the individual formal
distributors of \cref{def:formaldistributor}. A horizontal boundary is
a finite path of these distributors, including their intervening formal
categories. Each individual distributor also retains its carrier path
in the original bicomodule graph.
\end{definition}

Thus there are two levels of paths in \cref{eq:FDist}. The carrier of
one distributor is a bicomodule path in $\mathbb B(\V)$. A boundary
in $\FDist(\V)$ is a path of distributors. The latter carries the
intermediate formal-category actions encoded by \cref{eq:distbalance}.

\begin{theorem}[Conservative distributor completion]\label{thm:distcompletion}
The construction \cref{eq:FDist} has the properties of
\cref{thm:maindistributors}. More precisely, it induces natural bijections
\begin{equation}\label{eq:distconservative}
 \Ddist(\V)_{\mathsf F,\mathsf G}(\Gamma;N)
 \cong\FDist(\V)_{\mathsf F,\mathsf G}(\Gamma;N)
\end{equation}
for every distributor path $\Gamma$ and individual distributor $N$.
They preserve all distributor substitutions and reflect their equations.
Every many-sided cell has a finite expression in distributor multicells,
pasting, and inverses of intrinsic contextual comparisons.
\end{theorem}
\begin{proof}\leavevmode
\begin{enumerate}
\item \textbf{Apply the conservative completion.} Apply \cref{thm:completion,thm:globalembedding} to the virtual equipment
of \cref{cor:distequipment}. These theorems keep its object and vertical
categories, which are identified by \cref{prop:distvirtual}, and all
its designated horizontal arrows. They give
\cref{eq:distconservative}, its substitution compatibility, and the
poly-virtual equipment structure.

\item \textbf{Verify the finite presentation.} The finite-expression assertion is the
forest and localization presentation, together with the finite binding
construction used to globalize it. Each primitive distributor multicell
is itself a cell of $\mathbb B(\V)$ with the equations of
\cref{def:distcell}. Thus the whole construction is specified by the
existing finite-term calculus and the displayed equations.\qedhere
\end{enumerate}
\end{proof}

\begin{corollary}[Preservation of actual bicomodule distributors]\label{cor:atomicdistributors}
The formal categories, formal functors, and bicomodule-carried distributors
form a unital virtual double category with the multicells of
\cref{def:distcell}. Its inclusion in $\FDist(\V)$ preserves and
reflects all its single-output cells. In particular, for actual bicomodule carriers
$M,N$, its unary cells are exactly actual equivariant bicomodule maps
$\theta:M\to N$ satisfying \cref{eq:modulemap}.
\end{corollary}
\begin{proof}\leavevmode
\begin{enumerate}
\item \textbf{Close the atomic fragment under substitution.} The indicated restriction on horizontal arrows is closed under virtual
substitution, which does not form a new carrier. It contains the regular
distributor at every object.

\item \textbf{Apply distributor conservativity.} Its multicells are the same as those of
$\Ddist(\V)$ on these profiles, so \cref{eq:distconservative}
gives the assertion.

\item \textbf{Identify unary distributor maps.} For unary cells, \cref{thm:arbitraryV} identifies
the underlying square with an actual equivariant bicomodule map, and
\cref{eq:distexterior} become exactly
\cref{eq:modulemap}.\qedhere
\end{enumerate}
\end{proof}

\begin{proposition}[Functor companions and conjoints]\label{prop:distmates}
For a formal functor $\mathsf F=(f,F):\mathsf A\to\mathsf B$,
any local choice of the restrictions
\begin{equation}\label{eq:distcompanions}
 \mathsf F_*=U_{\mathsf B}[\mathsf F,1_{\mathsf B}],
 \qquad
 \mathsf F^*=U_{\mathsf B}[1_{\mathsf B},\mathsf F].
\end{equation}
gives companions and conjoints; their underlying carrier paths are
$f_*B$ and $Bf^*$, respectively.
In the completed calculus they have unit and counit
\[
 j_{\mathsf F}:\emp{\mathsf A}\cell
       (\mathsf F_*,\mathsf F^*),\qquad
 e_{\mathsf F}:(\mathsf F^*,\mathsf F_*)\cell\emp{\mathsf B},
\]
with the triangle identities. There are coherent bijections
\begin{equation}\label{eq:distmates}
 \FDist(\V)_{\mathsf F,\mathsf G}(\Gamma;\Delta)
 \cong
 \FDist(\V)_{\mathrm{glob}}
      (\Gamma;\mathsf F_*\Delta\mathsf G^*).
\end{equation}
They respect the defined tree substitutions and compositions of formal
functors.
\end{proposition}
\begin{proof}\leavevmode
\begin{enumerate}
\item \textbf{Construct the distributor restrictions.} The two restrictions in \cref{eq:distcompanions} exist by
\cref{thm:distrestriction}; \cref{eq:distrestrictioncarrier} computes
their carriers.

\item \textbf{Obtain coherent binding data.} Use these restriction cells and the unit comparison of
\cref{lem:distunit} in \cref{lem:binding}, now with input
$\Ddist(\V)$. That lemma constructs $j_{\mathsf F},e_{\mathsf F}$,
proves the triangle identities, and supplies the coherent comparisons
for composite functors.

\item \textbf{Construct the mate bijections.} \Cref{prop:mates} gives
\cref{eq:distmates}. Its inverse and its second mate description use
the single-wire cuts of that proposition, now on distributor wires.

\item \textbf{Verify composition and substitution compatibility.} The unit and counit composition equations prove compatibility with
vertical composition, and tree substitution compatibility follows by
substituting the same binding cuts into each tree.

\item \textbf{Normalize identity companions.} Identity companions
may be normalized to empty distributor paths as in
\cref{sec:global}.\qedhere
\end{enumerate}
\end{proof}

\begin{remark}[The second comparison class]\label{rem:secondcomparison}
The unit $\eta_A:\emp{\mathsf A}\cell U_{\mathsf A}$ is
contextually universal in the distributor calculus by
\cref{lem:distunit}. Its underlying cell
$\eta_A:\emp X\cell A$ in $\mathbb B(\V)$ need not be
invertible. For example, over the unit comonoid in $\Set$, take the
two-element group as a monoid. Its unit map from the singleton is not
an isomorphism; unary conservativity and the represented singleton
identify this with noninvertibility of the underlying unit cell in
$\mathbb B(\Set)$.

Forgetting distributor actions defines a virtual functor to
$\Kcal(\V)$, but this example shows that it does not in general
preserve intrinsic contextual comparisons. By
\cref{thm:functoriality}, it therefore need not extend across the
second localization with the same values on cells. The second completion
uses the distributor universal properties, including its regular units.
\end{remark}

The argument applies equally to any merge equipment, presented with paths
as horizontal arrows, together with a specified collection of its monoids:
take its underlying virtual equipment, form modules, restrict on those
objects, and apply the conservative completion. In the present application
the merge realization is already part of the construction of $\Q(\V)$,
and the chosen monoids are exactly its formal categories on actual
bicomodules.

\section{Representing distributor carriers and composition}\label{sec:distributorrecovery}

\begin{lemma}[Transport of a distributor carrier]\label{lem:distcarriertransport}
Let $M:\mathsf A\harr\mathsf B$ have carrier $H$, and let
$h:H\cell T$ be an invertible globular cell in $\mathbb B(\V)$.
There is a unique distributor structure on $T$ making $h$ an equivariant
isomorphism from $M$. Its actions are
\begin{equation}\label{eq:disttransport}
 p_T=h p_M(1_A,h^{-1}),\qquad
 q_T=h q_M(h^{-1},1_B).
\end{equation}
Transport also gives bijections on balanced distributor multicells and
preserves all their substitutions.
\end{lemma}
\begin{proof}\leavevmode
\begin{enumerate}
\item \textbf{Determine the transported actions.} The equations for $h$ to preserve the actions force
\cref{eq:disttransport}.

\item \textbf{Verify the action laws.} Substitution cancels each adjacent
$h^{-1}h$ in the associative and commuting-action equations, reducing
them to those for $M$. The unit equations reduce to
$h1_Hh^{-1}=1_T$. Thus the actions exist and satisfy all the laws.

\item \textbf{Transport multicells and their equations.} For a multicell, postcompose with the comparison on its output and
substitute the inverse comparisons on its inputs. Each exterior or
balance equation becomes the corresponding old equation after these
cancellations. The same argument covers \cref{eq:distnullary}.

\item \textbf{Verify invertibility and substitution compatibility.} Inverse transport uses the inverse comparisons, and substitution
compatibility follows from cancellation at every internal carrier.\qedhere
\end{enumerate}
\end{proof}

\begin{corollary}[Cotensor recovery for distributor carriers]\label{cor:distcotensor}
A distributor whose carrier path is stably represented by an actual
bicomodule is canonically isomorphic to a bicomodule-carried distributor.
When its action paths are also stably represented, its structures and
equivariant maps are exactly the cotensor bimodule structures and maps
of \cref{thm:bimodules}. The identifications respect changing bases and
balanced distributor multicells.
\end{corollary}
\begin{proof}\leavevmode
\begin{enumerate}
\item \textbf{Transport actions and maps to the representative.} Use the inverse cone comparison as $h$ in
\cref{lem:distcarriertransport}. Apply \cref{thm:bimodules} to the
transported actions, and \cref{lem:changebase} to exterior functors
and equivariant maps.

\item \textbf{Transport multicells.} For multicells, transport every carrier first.
The bijection of \cref{lem:distcarriertransport} respects their balance
and exterior equations.

\item \textbf{Identify maps from represented sources.} For a stably represented source path, unary full
faithfulness further identifies the transported cell with a map from the
corresponding cotensor; its equations are exactly equivariance and
balancing of that map.

\item \textbf{Prove independence of parentheses.} Cone uniqueness gives independence of the chosen
parentheses in each represented path.\qedhere
\end{enumerate}
\end{proof}

\begin{definition}[Represented distributor path]\label{def:representeddist}
A globular distributor multicell $c:\Gamma\cell P$ represents
$\Gamma$ when substitution gives bijections
\begin{equation}\label{eq:distuniversal}
 \Ddist(\V)(L,P,R;N)
 \longrightarrow\Ddist(\V)(L,\Gamma,R;N)
\end{equation}
for every compatible distributor context $L,R$ and every individual
distributor $N$. These are exactly the intrinsic comparisons of the
globular calculus of $\Ddist(\V)$. By globular detection, they also
satisfy the same property over arbitrary exterior formal functors.
\end{definition}

\begin{proposition}[Represented distributor substitution]\label{prop:distrepresented}
For a representing cell $c:\Gamma\cell P$, the completed merge
realization makes $c$ invertible. In its poly-virtual reduct,
\begin{equation}\label{eq:distconversion}
 a\longmapsto a(1_L,c,1_R):
 \FDist(\V)(L,P,R;\Delta)
 \longrightarrow\FDist(\V)(L,\Gamma,R;\Delta)
\end{equation}
is a bijection for every compatible distributor path $\Delta$ and
every exterior pair of formal functors. Its forward operation is a
single cut along the distributor $P$. Conversely an original
$c:\Gamma\cell P$ is invertible in the completion exactly when it
has \cref{eq:distuniversal}.
\end{proposition}
\begin{proof}\leavevmode
\begin{enumerate}
\item \textbf{Apply distributor detection.} Apply \cref{thm:detection,thm:globalembedding} to $\Ddist(\V)$.

\item \textbf{Obtain contextual factorization.} In the merge realization, contextual precomposition with an invertible
cell is a bijection.

\item \textbf{Retain the conversion in the reduct.} The forward map in \cref{eq:distconversion} is
a single-wire tree cut, and the reduct retains its values, all the cell
collections, and their equalities. It therefore retains the bijection.

\item \textbf{Prove the converse.} The converse is the detection theorem for this second comparison class.\qedhere
\end{enumerate}
\end{proof}

We give a construction of representing distributor cells. The following
definition describes a property of an existing coequalizer in the merge
realization; the definition of $\FDist(\V)$ is independent of its
existence.

\begin{definition}[Contextually preserved action coequalizer]\label{def:distcoequalizer}
Let $M:\mathsf A\harr\mathsf B$ and
$N:\mathsf B\harr\mathsf C$ have carriers $H_M,H_N$, where the
middle formal category has carrier $B$. An action coequalizer is a fork
in the relevant globular hom-category of $\mathbb B(\V)$,
\begin{equation}\label{eq:distcoequalizer}
 \begin{tikzcd}[column sep=3em]
 H_M B H_N
 \arrow[r,shift left=.7ex,"q_M\wh1"]
 \arrow[r,shift right=.7ex,"1\wh p_N"'] &
 H_M H_N\arrow[r,"c"] &T,
 \end{tikzcd}
\end{equation}
in which $c$ is a coequalizer. It is contextually preserved when, for
every compatible pair of bicomodule paths $L,R$, whiskering the entire
fork by $L$ and $R$ again gives a coequalizer in the corresponding
globular hom-category. The carrier $T$ is a horizontal path of
$\mathbb B(\V)$.
\end{definition}

\begin{theorem}[Relative tensor construction]\label{thm:distcoequalizer}
A contextually preserved action coequalizer
\cref{eq:distcoequalizer} gives a unique distributor structure
$P:\mathsf A\harr\mathsf C$ on $T$ making
$c:(M,N)\cell P$ a distributor multicell. This multicell represents
$(M,N)$ in $\Ddist(\V)$.
\end{theorem}
\begin{proof}\leavevmode
\begin{enumerate}
\item \textbf{Construct the exterior actions.} Let $A,C$ be the carriers of the two exterior formal categories. The
cell $c(p_M\wh1_{H_N})$ coequalizes the two action maps with $A$
adjoined on the left. To check this, commute the left $A$-action on
$M$ past its right $B$-action, and then use the equality defining
$c$. Preservation of the fork by left whiskering with $A$ therefore
gives a unique left action $p_P$ with
\begin{equation}\label{eq:distcoeqactions}
 p_P(1_A\wh c)=c(p_M\wh1_{H_N}),\qquad
 q_P(c\wh1_C)=c(1_{H_M}\wh q_N),
\end{equation}
where the second action is constructed by the right-hand version of the
same argument.

\item \textbf{Verify action associativity and units.} Every contextual instance of $c$ is an epimorphism, since it is a
coequalizer. Precompose the proposed left associativity equation with
$1_A\wh1_A\wh c$. Using \cref{eq:distcoeqactions}, it becomes
left associativity for $p_M$, postcomposed with $c$. Epimorphy gives
the desired equality. The left unit equation is tested with $c$ and
reduces to the unit equation for $p_M$. Right associativity and its
unit equation are proved by the symmetric calculation with $q_N$.

\item \textbf{Verify commuting actions.} The commuting-action equation is tested with
$1_A\wh c\wh1_C$; it reduces to the equality of the two actions
on the disjoint factors $H_M,H_N$, by interchange. Hence $P$ is a
distributor.

\item \textbf{Identify the distributor comparison.} The two equations in \eqref{eq:distcoeqactions} are its exterior
equivariance equations for $c$, and the coequalizer equation is its
interior balance. They also prove uniqueness of the distributor actions.

\item \textbf{Factor a balanced multicell.} For contextual universality, let
$\alpha:(L,M,N,R)\cell Z$ be a globular balanced distributor multicell.
Its balance at the middle formal category says exactly that the
underlying cell coequalizes the whiskered pair in
\cref{eq:distcoequalizer}. The preserved coequalizer therefore gives a
unique underlying factor $\bar\alpha:(L,P,R)\cell Z$.

\item \textbf{Reflect balance and exterior equations.} To verify
any remaining balance or exterior equation for $\bar\alpha$,
precompose with the corresponding contextual instance of $c$, also
including the extra category carrier in that equation. The result is
the corresponding equation for $\alpha$, using
\cref{eq:distcoeqactions} at the two interfaces adjacent to $P$.
That contextual instance is epic, so the equation holds.

\item \textbf{Prove the contextual bijection.} Conversely,
substitution of $c$ into a distributor multicell is a distributor
multicell by \cref{prop:distvirtual}. These two operations are inverse
by the coequalizer property. This proves \cref{eq:distuniversal}.

\item \textbf{Allow changing bases.} Globular detection in the virtual equipment $\Ddist(\V)$ gives all
changing-base instances. No preservation of arbitrary coequalizers or
ambient exactness property is used: only the stated contextual preservation
of the particular fork~\eqref{eq:distcoequalizer}.\qedhere
\end{enumerate}
\end{proof}

When these representations exist for successive distributor pairs, their
iterated composites have the expected coherent associators and unitors.
Indeed substitution of intrinsic comparisons is intrinsic by
\cref{lem:universalclosure}. Any two representatives of a path therefore
have a unique invertible distributor map carrying one representing cell
to the other. Every route around a pentagon or triangle carries the same
representing cell to the same cell, so uniqueness proves equality. The
same factorization property defines maps of represented composites over
formal functors and proves their functoriality.

For bicomodule-carried distributors in a represented cotensor setting,
the two objects before $T$ in \cref{eq:distcoequalizer} are represented
by $H_M\square B\square H_N$ and $H_M\square H_N$, with the
indicated intermediate comonoid bases. Thus the pair is the familiar
pair of module actions used to form a relative tensor product. Its
contextual coequalizer property is a separate representation property
from the stable equalizers used to construct those cotensors. The next
section verifies both properties in a nonregular example.

\section{Examples and the role of regularity}\label{sec:examples}

The unit and Cartesian bases recover familiar categories. Finite
group-like bases in $\Ab$ recover enriched categories and distributors
despite ambient nonregularity. In \cref{sec:unrepresented}, instability
instead occurs within an action boundary on one actual bicomodule.

\begin{proposition}[The unit comonoid]\label{prop:unitbase}
For every monoidal $\V$, formal internal categories over $I$ are
exactly ordinary monoids in $\V$.
\end{proposition}
\begin{proof}\leavevmode
\begin{enumerate}
\item \textbf{Identify the unit-base coactions.} The counit law forces either $I$-coaction to be the inverse unitor.

\item \textbf{Compute stable representatives.} The two maps defining balance over $I$ are therefore equal after
the canonical unit identifications. Their equalizer is the identity
on $M\otimes N$, which is preserved by every tensor functor.
Thus every finite path over $I$ is stably represented by its ordinary
tensor product, and the nullary representative is $I$.

\item \textbf{Recover ordinary monoids.} \Cref{thm:aguiar,thm:globalaguiar} identify the resulting structures
and fixed-base maps with ordinary monoids and their homomorphisms.\qedhere
\end{enumerate}
\end{proof}

\begin{proposition}[Identity category]\label{prop:identitycategory}
Every comonoid $C$ in every monoidal $\V$ has a formal internal
category with carrier its regular bicomodule $C$.
\end{proposition}
\begin{proof}\leavevmode
\begin{enumerate}
\item \textbf{Transport the identity monoid.} By \cref{lem:bicomodunit}, the comparison
$d_C:C\cell\emp C$ is invertible in the merge realization, with
inverse $\tau_0$. Transport the identity monoid on the empty path
along this isomorphism. Explicitly,
\[
 \eta=\tau_0,\qquad
 \mu=\tau_0(d_C\wh d_C):(C,C)\cell C.
\]

\item \textbf{Verify the formal monoid laws.} Substitution cancels the adjacent pairs $d_C\tau_0$, proving
associativity and the two unit laws. They remain true in the reduct
by \cref{thm:reduct}.

\item \textbf{Identify the diagonal comparison.} Equivalently, the multiplication is the
inverse comparison for $\delta_C:C\cell(C,C)$: the unit cone
calculation in \cref{lem:cotensorcoherence} identifies
$C\square_C C$ with $C$ in every monoidal category.\qedhere
\end{enumerate}
\end{proof}

\begin{proposition}[Cartesian bases]
For a cartesian monoidal category with finite limits, formal internal
categories are ordinary internal categories, with their usual internal
functors.
\end{proposition}
\begin{proof}\leavevmode
\begin{enumerate}
\item \textbf{Identify bicomodules with internal graphs.} Every object has its unique diagonal comonoid. Counitality forces a
left coaction $M\to C\times M$ to have second component $1_M$,
so it is equivalent to a map $M\to C$; the right coaction likewise
amounts to a map $M\to D$. Coassociativity and commuting coactions
then follow from the product identities. Thus a bicomodule is an
internal graph with its two endpoint maps.

\item \textbf{Identify cotensors with stable pullbacks.} The balance equalizer of
$M\times N\rightrightarrows M\times D\times N$ is exactly
the pullback $M\times_D N$. Product with any object preserves
this equalizer: a map into the product is a pair of maps and factors
exactly when its second component does. The same holds on both
sides, so the cone is stable.

\item \textbf{Recover internal categories and functors.} The maps and equations of
\cref{thm:aguiar,thm:globalaguiar} now say precisely that composable
arrows compose associatively, identity arrows are units, and functors
preserve both operations.\qedhere
\end{enumerate}
\end{proof}

\begin{proposition}[A nonregular ambient category]\label{prop:Ab}
The monoidal category $(\Ab,\otimes_{\mathbb Z},\mathbb Z)$ is not
regular in Aguiar's sense. For
\[
 C=\bigoplus_{i=1}^{r}\mathbb Z e_i,\qquad
 \delta_C(e_i)=e_i\otimes e_i,\quad \eps_C(e_i)=1,
\]
all bicomodule paths over $C$ nevertheless have stable cotensors, with
\begin{equation}\label{eq:matrixcotensor}
 (M\square_C N)_{ik}=\bigoplus_{j=1}^{r}M_{ij}\otimes N_{jk}.
\end{equation}
Formal categories over $C$ are categories enriched in $\Ab$ on the
fixed object set $\{1,\ldots,r\}$.
\end{proposition}
\begin{proof}\leavevmode
\begin{enumerate}
\item \textbf{Exhibit failure of ambient regularity.} The coreflexive pair $r,s:\mathbb Z\rightrightarrows\mathbb Z\oplus
\mathbb Z$ given by
\[
 r(n)=(n,2n),\qquad s(n)=(n,0)
\]
has equalizer zero; projection to the first coordinate is a common
retraction. After tensoring with $\mathbb Z/2$, the two maps are
equal, so their equalizer is $\mathbb Z/2$. Hence tensoring does not
preserve this equalizer.

\item \textbf{Decompose bicomodules into matrix entries.} For the specified comonoid write
\[
 \lambda_M(x)=\sum_i e_i\otimes p_i(x),\qquad
 \rho_M(x)=\sum_j q_j(x)\otimes e_j.
\]
Coassociativity gives $p_ip_k=\delta_{ik}p_i$ and
$q_jq_l=\delta_{jl}q_j$; counitality gives
$\sum_i p_i=1=\sum_j q_j$. Commuting coactions gives
$p_iq_j=q_jp_i$. Thus
$M=\bigoplus_{i,j}M_{ij}$, where $M_{ij}=\operatorname{im}(p_iq_j)$.
The converse constructs the coactions from any such decomposition.

\item \textbf{Compute the matrix cotensor.} Decompose $M\otimes N$ into the summands
$M_{ij}\otimes N_{kl}$. The two balance maps insert $e_j$ and
$e_k$, respectively. If $j=k$ they agree. If $j\ne k$, projection
to the inserted $e_j$ coordinate takes their difference to the
identity on this summand. Therefore the equalizer consists exactly
of the summands with $j=k$, giving \cref{eq:matrixcotensor}.

\item \textbf{Prove stability under tensoring.} This argument remains valid after arbitrary tensoring on either side:
finite direct sums are preserved, and the coordinate calculation still
gives identity or zero on each summand. It proves preservation of the
equalizer fork, not just preservation of its inclusion. In particular, this
stability is a coordinate calculation after tensoring and does not use
flatness of any of the groups involved.

\item \textbf{Recover enriched category composition.} For $A:C\harr C$, a multiplication consists of maps
$A_{ij}\otimes A_{jk}\to A_{ik}$. A bicomodule map $C\to A$
consists of maps $\mathbb Z\to A_{ii}$. The cotensor associator is
the canonical associator for the direct sums of triple tensor products,
by its cone characterization. Thus the monoid equations are precisely
bilinear category composition and its identity equations. Apply
\cref{thm:aguiar} to these stable cones.\qedhere
\end{enumerate}
\end{proof}

\begin{remark}
The preceding identification includes the variation of the finite
object set. A comonoid map
$\bigoplus_i\mathbb Z e_i\to\bigoplus_j\mathbb Z e_j$ must send
each $e_i$ to a group-like element. Writing such an element as
$\sum_j a_je_j$, its equations give $a_j^2=a_j$, $a_ja_k=0$ for
$j\ne k$, and $\sum_j a_j=1$. Over $\mathbb Z$ exactly one
coefficient is $1$ and the rest are $0$. Thus these comonoid maps are
exactly functions of the indexing sets, and \cref{thm:globalaguiar}
gives the enriched functors.
\end{remark}

In this nonregular example, every cotensor relevant to the indicated
finite bases is still represented stably. Its role is to show why a global
regularity condition is stronger than the requirements of these category
structures. The definition in \cref{def:Q} also applies to paths without
such representatives, using the same formal cell calculus.

\begin{proposition}[Finite-object enriched distributors]\label{prop:Abdistributors}
Apply the distributor construction to the part of $\mathbb B(\Ab)$
whose objects are finite group-like comonoids
$C_S=\bigoplus_{s\in S}\mathbb Z e_s$ and whose carrier paths have
all vertices of this form. Its objects and vertical arrows are
finite-object $\Ab$-enriched categories and enriched functors. Every
distributor is isomorphic to one carried by an actual bicomodule, and
these distributors are exactly enriched distributors. Every finite
distributor path is represented by its relative tensor product.
Consequently, after replacing represented boundaries, the completed
many-sided cells are exactly the equivariant maps between these relative
tensor products.
\end{proposition}
\begin{proof}\leavevmode
\begin{enumerate}
\item \textbf{Restrict to finite group-like bases.} Let $\mathbb B_{\mathrm{fin}}$ be the specified part of
$\mathbb B(\Ab)$. The companion and conjoint of a comonoid map
$C_S\to C_T$ are bicomodules with these finite bases, so the
restriction paths remain in $\mathbb B_{\mathrm{fin}}$. Units remain
there as well. Thus its path-carrier calculus is a virtual equipment
by the proof of \cref{lem:carrierequipment}.

\item \textbf{Represent carrier paths by matrices.} \Cref{prop:Ab} and its proof, with rectangular indexing sets, identify
a bicomodule $C_S\harr C_T$ with a matrix of abelian groups
$(M_{ij})_{i\in S,j\in T}$ and identify cotensor composition with
matrix multiplication using direct sums and tensor products. All carrier
paths in $\mathbb B_{\mathrm{fin}}$ therefore have stable atomic
representatives.

\item \textbf{Identify enriched distributor data.} \Cref{lem:distcarriertransport} replaces every
distributor carrier by such a representative. Its actions are maps
\[
 A_{ii'}\otimes M_{i'j}\longrightarrow M_{ij},\qquad
 M_{ij}\otimes B_{jj'}\longrightarrow M_{ij'},
\]
with the associative, unital, and commuting equations of an enriched
distributor. The exterior equations on maps give the standard equivariant
maps over enriched functors. Balanced multicells have the same
identification by \cref{def:distcell} and cotensor representation.

\item \textbf{Construct the balanced tensor quotient.} For distributors $M:\mathsf A\harr\mathsf B$ and
$N:\mathsf B\harr\mathsf C$, define
\begin{equation}\label{eq:Abdistributorproduct}
 (M\otimes_{\mathsf B}N)_{ik}
 =\frac{\displaystyle\bigoplus_j M_{ij}\otimes N_{jk}}
 {\bigl\langle (mb)\otimes n-m\otimes(bn)\bigr\rangle}.
\end{equation}
Here the subgroup is generated by all indicated differences, with
$m\in M_{ij}$, $b\in B_{jj'}$, and $n\in N_{j'k}$; the two terms
belong to the $j'$ and $j$ summands, respectively. For fixed exterior indices $i,k$, put
\[
X=\bigoplus_{j,j'}M_{ij}\otimes B_{jj'}\otimes N_{j'k},\quad
Y=\bigoplus_j M_{ij}\otimes N_{jk},\quad
Q=(M\otimes_{\mathsf B}N)_{ik}.
\]
The componentwise coequalizer of the two middle actions has the
balanced-map factorization
\[
\begin{tikzcd}[column sep=3em]
X \arrow[r,shift left=.7ex,"q_M\otimes1"]
  \arrow[r,shift right=.7ex,"1\otimes p_N"'] &
Y \arrow[r,two heads,"\pi"] \arrow[dr,"\alpha"'] &
Q \arrow[d,dashed,"\exists!\,\bar\alpha"]\\
&& Z.
\end{tikzcd}
\]

\item \textbf{Verify contextual preservation of the quotient.} Tensor product of abelian groups preserves coequalizers in either variable, and finite
direct sums preserve them as well. Hence matrix cotensor composition
preserves this fork under every compatible matrix context. Replacing
path carriers by their representatives transfers this contextual
coequalizer property to $\mathbb B_{\mathrm{fin}}$.

\item \textbf{Represent distributor pairs and finite paths.} \Cref{thm:distcoequalizer} equips the quotient in
\eqref{eq:Abdistributorproduct} with its exterior actions and proves
that it represents the pair as distributors. The regular distributor
represents the empty path by \cref{lem:distunit}, so induction gives
all finite paths. Coherence and functoriality follow from the unique
representing comparisons.

\item \textbf{Identify cells in the second completion.} In the second completion, replace both
boundaries by these representatives and use
\cref{eq:distconservative} to identify the resulting unary square
with its ordinary equivariant distributor map. This proves the final
assertion, including arbitrary exterior enriched functors.\qedhere
\end{enumerate}
\end{proof}

In this example a functor $\mathsf F:\mathsf A\to\mathsf B$ with
object function $f:S\to T$ has the familiar represented companion and
conjoint
\[
 (\mathsf F_*)_{ij}=B_{f(i),j},\qquad
 (\mathsf F^*)_{ji}=B_{j,f(i)}.
\]
Indeed these are the matrix cotensors of the carrier paths $f_*B$ and
$Bf^*$ in \cref{eq:distcompanions}; their actions are exactly the
restrictions constructed in \cref{thm:distrestriction}. This comparison
is made in the indicated finite-base construction, where every carrier
path is represented. The global construction of \cref{eq:FDist} also
retains distributor carriers that pass through other comonoid bases.

\section{A distributor with an unrepresented action boundary}
\label{sec:unrepresented}

The construction has two steps. First form a category whose own
composition paths are stable. Then restrict a stable distributor action
along a formal functor. The resulting action remains a well-defined
cell, while its input path loses contextual representability.

\subsection{The category and the functor}

\begin{definition}[Triangular data]\label{strong:def:triangular-data}
Let $D=\Z g\oplus\Z x$ with
\[
\delta(g)=g\otimes g,\qquad
\delta(x)=g\otimes x+x\otimes g,\qquad
\eps(g)=1,\quad\eps(x)=0.
\]
Set $E=\Z e\oplus D$, where $e$ is group-like. Let
$H=\Z a\oplus\Z b$ be the $E$--$E$ bicomodule
\begin{equation}\label{strong:eq:H}
\lambda_H(h)=e\otimes h,\qquad
\rho_H(a)=a\otimes g,\qquad
\rho_H(b)=b\otimes g+2a\otimes x.
\end{equation}
Also let $J=\Z t$ with $\lambda_J(t)=e\otimes t$ and
$\rho_J(t)=t\otimes g$, and define $\varphi:H\to J$ by
$\varphi(a)=0$, $\varphi(b)=t$.
\end{definition}

\begin{lemma}[Validity of the triangular data]\label{strong:lem:triangular-data}
The displayed formulas make $E$ a comonoid, $H,J$ bicomodules over
$E$, and $\varphi$ an equivariant bicomodule map.
\end{lemma}
\begin{proof}\leavevmode
\begin{enumerate}
\item \textbf{Verify the comonoid laws.} The comonoid laws for $E$ hold on $e,g,x$: the two iterated coproducts
of $x$ are the sum of the three tensors with one $x$ and two $g$'s,
and the counit laws follow from $\eps(g)=1$, $\eps(x)=0$.

\item \textbf{Verify the bicomodule laws.} For the right coaction on $H$, write
$\rho_H(h)=h\otimes g+d(h)\otimes x$ with $d(a)=0$, $d(b)=2a$.
The equation $d^2=0$ proves coassociativity, the first square below.
Counitality and commuting coactions follow from the formulas; the
laws for $J$ use that $e,g$ are group-like.

\item \textbf{Verify equivariance of the map.} Right equivariance of $\varphi$ is the second square, since
$\varphi d=0$; left equivariance follows from the common support $e$:
\[
\begin{tikzcd}[column sep=2.1em]
H \arrow[r,"\rho_H"] \arrow[d,"\rho_H"'] &
H\otimes E \arrow[d,"1\otimes\delta_E"]\\
H\otimes E \arrow[r,"\rho_H\otimes1"'] & H\otimes E\otimes E
\end{tikzcd}
\qquad
\begin{tikzcd}[column sep=2.1em]
H \arrow[r,"\varphi"] \arrow[d,"\rho_H"'] &
J \arrow[d,"\rho_J"]\\
H\otimes E \arrow[r,"\varphi\otimes1"'] & J\otimes E.
\end{tikzcd}\qedhere
\]
\end{enumerate}
\end{proof}

\begin{lemma}[Triangular categories]\label{strong:lem:triangular}
For $K=H$ or $J$, the carrier $A_K=E\oplus K$ defines a formal category,
with unit the inclusion of $E$ and multiplication the fold
\[
m_K:E\oplus K\oplus K\longrightarrow E\oplus K,\qquad
(c,k_1,k_2)\longmapsto(c,k_1+k_2).
\]
Every finite composition path is stably represented; for $n\geq1$,
\[
A_K^{\square_E n}\cong E\oplus K^{\oplus n}.
\]
\end{lemma}
\begin{proof}\leavevmode
\begin{enumerate}
\item \textbf{Compute the stable binary blocks.} Expand $(E\oplus K)\otimes(E\oplus K)$ into its four blocks.
The $E,E$, $E,K$, and $K,E$ blocks have the stable unit cotensors
$E$, $K$, and $K$. The $K,K$ block has stable equalizer zero:
projection of the inserted $E$ factor onto its $e$ coordinate takes
the two balance maps to zero and the identity, respectively. Thus every
equalizing map is zero, also after tensoring.

\item \textbf{Represent every finite word.} For a longer word, remove regular $E$ factors using their stable unit
cones. A block with no $K$ contributes $E$, one $K$ contributes $K$,
and two or more $K$'s contribute zero. This proves the asserted
representatives and stability.

\item \textbf{Prove associativity of the fold.} Associativity is the fold square
\[
\begin{tikzcd}[column sep=4em]
E\oplus K^{\oplus3} \arrow[r,"\nabla_{12}"]
  \arrow[d,"\nabla_{23}"'] &
E\oplus K^{\oplus2} \arrow[d,"m_K"]\\
E\oplus K^{\oplus2} \arrow[r,"m_K"'] & E\oplus K.
\end{tikzcd}
\]
Here $\nabla_{12}$ and $\nabla_{23}$ fold the indicated pair of $K$
summands; both routes add all three.

\item \textbf{Verify units and transport the laws.} The unit
equations are the summand inclusions. \Cref{thm:aguiar}
converts these into formal category laws.\qedhere
\end{enumerate}
\end{proof}

\begin{corollary}[The restriction functor]\label{strong:cor:triangular-functor}
Write $A=E\oplus H$ and $B=E\oplus J$, with the resulting categories
$\mathsf A,\mathsf B$. The map
\[
F=1_E\oplus\varphi:A\longrightarrow B
\]
defines a formal functor $\mathsf F=(1_E,F):\mathsf A\to\mathsf B$.
\end{corollary}
\begin{proof}\leavevmode
\begin{enumerate}
\item \textbf{Verify equivariance.} By \cref{strong:lem:triangular-data}, $F$ is equivariant.

\item \textbf{Verify the functor equations.} It preserves
the unit and the folds of \cref{strong:lem:triangular} through
\[
\begin{tikzcd}[column sep=3.8em]
E\oplus H^{\oplus2} \arrow[r,"m_H"]
  \arrow[d,"1_E\oplus\varphi^{\oplus2}"'] &
A \arrow[d,"F"]\\
E\oplus J^{\oplus2} \arrow[r,"m_J"'] & B
\end{tikzcd}
\qquad
\begin{tikzcd}
E \arrow[r,hook,"u_A"] \arrow[dr,"u_B"'] &
A \arrow[d,"F"]\\
& B.
\end{tikzcd}
\]
\Cref{thm:globalaguiar} gives the formal functor laws.

\item \textbf{Record the carrier ranks.} The carriers $A,B,E$ have ranks five, four, and three, respectively.\qedhere
\end{enumerate}
\end{proof}

\subsection{An actual bicomodule with a nonzero formal action}

\begin{definition}[The distributor carrier]\label{strong:def:M}
Put $N=\Z/4$ and $M=Nu\oplus Nv$, with
\[
\lambda_M(u)=e\otimes u,\qquad
\lambda_M(v)=g\otimes v,
\]
and the right coaction of the unit comonoid $\Z$.
Let $\Icat$ denote the identity formal category at $\Z$.
\end{definition}

\begin{proposition}[The stable action]\label{strong:prop:stable-action}
The carrier $M$ admits a $\mathsf B$--$\Icat$ distributor structure.
Its action paths $(B,M)$ and $(B,B,M)$ are stably represented by
$M\oplus N$ and $M\oplus N\oplus N$, respectively. Its left action
is induced by the actual equivariant map
\[
p_{B,0}:M\oplus N\longrightarrow M,\qquad
((ru+sv),w)\longmapsto(r+w)u+sv.
\]
\end{proposition}
\begin{proof}\leavevmode
\begin{enumerate}
\item \textbf{Compute the stable action cotensor.} For $(B,M)$, the regular $E$ block gives $M$;
$J\otimes Nu$ has mismatched supports and gives zero; and both
coactions on $J\otimes Nv$ insert $g$, giving all of $N$.
These computations persist under arbitrary tensoring.

\item \textbf{Check action equivariance.} The extra $N$ summand and its image $Nu$ both have left support $e$,
so this is equivariant.

\item \textbf{Verify associativity.} The triple $(B,B,M)$ is stably represented by
$M\oplus N\oplus N$, and the action square becomes
\[
\begin{tikzcd}[column sep=4em]
M\oplus N\oplus N \arrow[r,"1_M\oplus+"] \arrow[d,"\sigma"'] &
M\oplus N \arrow[d,"p_{B,0}"]\\
M\oplus N \arrow[r,"p_{B,0}"'] & M.
\end{tikzcd}
\]
Here $\sigma(m,w_1,w_2)=(p_{B,0}(m,w_2),w_1)$; both routes send
$((ru+sv),w_1,w_2)$ to $(r+w_1+w_2)u+sv$.

\item \textbf{Verify the action unit.} The unit selects $M$.
Hence this gives an associative, unital formal action $p_B$.

\item \textbf{Add the commuting right action.} The right action $q:(M,\Z)\cell M$ is the unit comparison and commutes
with $p_B$ by its coherence.\qedhere
\end{enumerate}
\end{proof}

\begin{proposition}[Restriction on the actual carrier]\label{strong:prop:restricted-action}
The same bicomodule $M$ admits a $\mathsf A$--$\Icat$ distributor
structure, obtained by restriction along
$\mathsf F$ of \cref{strong:cor:triangular-functor}. Its left action
is defined by the outer route
\begin{equation}\label{strong:eq:pA}
\begin{tikzcd}[column sep=4em]
(A,M) \arrow[r,"{(F,1_M)}"] \arrow[d,"p_A"'] &
(B,M) \arrow[d,"\tau_{B,M}","\sim"']\\
M & M\oplus N \arrow[l,"p_{B,0}"]
\end{tikzcd}
\end{equation}
where $\tau_{B,M}$ is the stable comparison already constructed.
Equivalently, $p_A=p_B(F,1_M)$, a tree substitution in $\Q(\Ab)$.
Its right action is the unit action $q$ of
\cref{strong:prop:stable-action}.
\end{proposition}
\begin{proof}\leavevmode
\begin{enumerate}
\item \textbf{Identify the restricted carrier.} Since the base map of $\mathsf F$ is $1_E$, restriction retains the
actual carrier $M$.

\item \textbf{Verify restricted associativity.} Associativity follows by pasting
\[
\begin{tikzcd}[column sep=2.5em,row sep=2.6em]
(A,A,M) \arrow[r,"{(\mu_A,1)}"] \arrow[d,"{(F,F,1)}"'] &
(A,M) \arrow[r,"p_A"] \arrow[d,"{(F,1)}"] &
M \arrow[d,equal]\\
(B,B,M) \arrow[r,"{(\mu_B,1)}"] \arrow[d,"{(1,p_B)}"'] &
(B,M) \arrow[r,"p_B"] \arrow[d,"p_B"] &
M \arrow[d,equal]\\
(B,M) \arrow[r,"p_B"'] & M \arrow[r,equal] & M.
\end{tikzcd}
\]
The two routes around the outside are the two $A$-action composites,
after using the definition of $p_A$ on the left route.

\item \textbf{Verify the restricted unit law.} Unitality is the left triangle in
\[
\begin{tikzcd}[column sep=3em]
M \arrow[r,"{(\eta_A,1)}"] \arrow[dr,equal] &
(A,M) \arrow[r,"{(F,1)}"] \arrow[d,"p_A" description] &
(B,M) \arrow[dl,"p_B"]\\
& M &
\end{tikzcd}
\]
whose outer triangle is the functor unit law followed by the $B$-action
unit law.

\item \textbf{Verify commuting actions.} The commuting-action equation follows by substituting $F$ into that for
$p_B,q$. Thus $(M,p_A,q)$ is a formal distributor
$\mathsf A\harr\Icat$ on one actual bicomodule. The construction of
$p_A$ uses the inverse comparison for the stable path $(B,M)$.\qedhere
\end{enumerate}
\end{proof}

\subsection{The unstable cone and the contextual obstruction}

\begin{lemma}[The unstable action cotensor]\label{strong:lem:unstable-cotensor}
The ordinary equalizer $T=A\square_E M$ exists and carries an
equivariant $E$--$\Z$ bicomodule cone. After tensoring with $R=\Z/2$,
the comparison from this cone to the new equalizer is neither monic
nor epic. In particular, the action cotensor is not stable.
\end{lemma}
\begin{proof}\leavevmode
\begin{enumerate}
\item \textbf{Compute the ordinary cotensor.} The regular block of $T$ is
$M$. The $H\otimes Nu$ block is zero by its supports, while an element
\[
a\otimes\alpha v+b\otimes\beta v\in H\otimes Nv
\]
is balanced exactly when $2\beta=0$ in $N$. Indeed its balance
obstruction is $2a\otimes x\otimes\beta v$. Consequently
\begin{equation}\label{strong:eq:badcotensor}
T\cong M\oplus Na\oplus N[2]b
\cong(\Z/4)^3\oplus\Z/2.
\end{equation}

\item \textbf{Describe its equivariant cone.} Writing $z$ for the last, order-two generator, its cone
$i:T\hookrightarrow A\otimes M$ is
\begin{equation}\label{strong:eq:badcone}
u\longmapsto e\otimes u,\quad
v\longmapsto g\otimes v,\quad
a\longmapsto a\otimes v,\quad
z\longmapsto2b\otimes v.
\end{equation}
Give $u,a,z$ left support $e$, and $v$ left support $g$; this makes
$T$ an $E$--$\Z$ bicomodule and $i$ an equivariant cone.

\item \textbf{Exhibit a kernel after tensoring.} After tensoring with $R=\Z/2$, the nonzero class
$z\otimes\overline1$ maps to zero.

\item \textbf{Exhibit a new balanced element.} Meanwhile
\begin{equation}\label{strong:eq:newbalanced}
b\otimes v\otimes\overline1
\end{equation}
becomes balanced: its obstruction is a multiple of $2$. This element
cannot lie in the image of $i\otimes1_R$, whose $b$ coordinate is zero.

\item \textbf{Compute the canonical comparison.} Thus the tensored cone fails both injectivity and surjectivity onto
the new equalizer $T_R^{\mathrm{bal}}$. The comparison is
\[
\begin{tikzcd}[column sep=4em]
T\otimes R \arrow[r,"i\otimes1_R"] \arrow[d,"\kappa"'] &
A\otimes M\otimes R \arrow[d,equal]\\
T_R^{\mathrm{bal}} \arrow[r,hook,"i_R"'] & A\otimes M\otimes R.
\end{tikzcd}
\]
Both objects on the left have underlying group $(\Z/2)^4$.
In the source basis $(u,v,a,z)$ and the target basis given by
$e\otimes u,g\otimes v,a\otimes v,b\otimes v$ after reduction,
$\kappa$ has matrix $\operatorname{diag}(1,1,1,0)$. The failure is in
this canonical cone, whose image has dimension three.\qedhere
\end{enumerate}
\end{proof}

\begin{theorem}[An unrepresented action boundary]\label{strong:thm:strong}\label{strong:prop:strong}
The action boundary $(A,M)$ has no atomic contextual representative.
Its formal action operates nontrivially on the new balanced element
\eqref{strong:eq:newbalanced}.
\end{theorem}
\begin{proof}\leavevmode
\begin{enumerate}
\item \textbf{Represent the empty-context functor.} In the original globular calculus $P$ of $\Ccal(\Ab)$, the cone $i$
is a cell $(A,M)\cell T$. It represents the empty-context balanced-map
functor: factor a balanced map through its ordinary equalizer, and
test equivariance of the factor after $1_E\otimes i$. This map is
monic because $E$ is free. The right coaction is the unit coaction.

\item \textbf{Construct the contextual obstruction.} Let $W=\Z/2$ have left $E$-coaction $w\mapsto e\otimes w$ and right
unit coaction. With $R$ regarded as a $\Z$--$\Z$ bicomodule, the map
\[
\beta:W\longrightarrow A\otimes M\otimes R,\qquad
\overline1\longmapsto b\otimes v\otimes\overline1
\]
satisfies both interior balances and exterior equivariance. It is a
cell in $P(A,M,R;W)$ for which the factorization triangle fails:
\[
\begin{tikzcd}[column sep=4.4em,row sep=2.6em]
W \arrow[r,"\beta"] \arrow[d,dashed,"\nexists\,h"'] &
A\otimes M\otimes R\\
T\otimes R \arrow[ur,"i\otimes1_R"'] &
\end{tikzcd}
\]
The dashed arrow denotes the nonexistent factor: the $b$ coordinate
of $\beta(\overline1)$ is nonzero, while that of every element in the
image of $i\otimes1_R$ vanishes.
Hence $i$ fails the contextual test of \cref{def:universal}.

\item \textbf{Exclude every alternative atomic representative.} Every other contextual representative would first represent the same
empty-context functor, so its cone would be $i$ composed with a
bicomodule isomorphism:
\[
\begin{tikzcd}
S \arrow[r,"\theta","\sim"'] \arrow[dr,"j"'] &
T \arrow[d,"i"]\\
& A\otimes M.
\end{tikzcd}
\]
Here $j$ is the proposed representing cone. Tensoring $\theta$ with $R$ cannot
change this failure of factorization. No alternative atom can therefore
represent the path. By \cref{thm:detection}, $(A,M)$ is not
isomorphic to any atom in the globular localization, or in its reversed
merge realization.

\item \textbf{Evaluate the formal action.} In the forward realization, $\beta$ is the primitive cell
$W\cell(A,M,R)$. Applying $F$ to its first output replaces $b$ by $t$.
The stable comparison for $(B,M)$ followed by $p_B$ then sends
$t\otimes v\otimes\overline1$ to $u\otimes\overline1$. Thus, in the
merge realization, the entire evaluation is
\[
\begin{tikzcd}[column sep=2.8em,row sep=3em]
W \arrow[r,"\beta"] \arrow[d,"\gamma"'] &
(A,M,R) \arrow[r,"{(F,1,1)}"] \arrow[d,"p_A\wh1_R"] &
(B,M,R) \arrow[d,"\tau_{B,M}\wh1_R","\sim"']\\
(M,R) \arrow[r,equal] &
(M,R) &
(M\oplus N,R) \arrow[l,"p_{B,0}\wh1_R"]
\end{tikzcd}
\]

\item \textbf{Detect the nonzero result.} Here $(M,R)$ is stably represented by $M\otimes R$.
Under this comparison, unary conservativity identifies $\gamma$ with
the actual map $\overline1\mapsto u\otimes\overline1\neq0$.
The reduct retains the resulting cells and equality.\qedhere
\end{enumerate}
\end{proof}

\begin{corollary}[The action on the ordinary cotensor]\label{strong:cor:ordinary-action}
The map $p_T:T\to M$ induced by the formal action is nonzero on the
order-two summand of $T$: it sends its generator $z$ to $2u$.
\end{corollary}
\begin{proof}\leavevmode
\begin{enumerate}
\item \textbf{Identify the ordinary cotensor action.} The primitive cone gives a triangle in the merge realization
\[
\begin{tikzcd}
T \arrow[r,"\iota"] \arrow[dr,"p_T"'] &
(A,M) \arrow[d,"p_A"]\\
& M,
\end{tikzcd}
\]
where unary conservativity identifies $p_T$ with the actual map
\[
p_T:T\longrightarrow M,\qquad
(ru+sv,\alpha,\beta)\longmapsto(r+\beta)u+sv,
\qquad \beta\in N[2],
\]
so $p_T(z)=2u\neq0$.

\item \textbf{Exclude contextual evaluation through that cotensor.} Nevertheless, the contextual action above cannot
be evaluated by first factoring through $T\otimes R$: that
factorization does not exist by \cref{strong:thm:strong}.\qedhere
\end{enumerate}
\end{proof}

\begin{corollary}[Restriction does not preserve action stability]\label{strong:cor:failure}
Stability of action cotensors need not survive restriction along a
formal functor, even when its base map is the identity and both formal
categories have stable finite composition paths. This failure occurs
for a distributor on one actual bicomodule, whose restricted action
boundary has no contextual atomic representative.
\end{corollary}
\begin{proof}\leavevmode
\begin{enumerate}
\item \textbf{Choose the functor and stable category structures.} Use the functor $\mathsf F:\mathsf A\to\mathsf B$ of
\cref{strong:cor:triangular-functor}. \Cref{strong:lem:triangular} gives
stable composition paths for both categories.

\item \textbf{Restrict the stable action.} The $\mathsf B$-action
of \cref{strong:prop:stable-action} has a stable action cotensor;
its restriction is the $\mathsf A$-action of
\cref{strong:prop:restricted-action}.

\item \textbf{Identify the failure of representation.} \Cref{strong:lem:unstable-cotensor}
and \cref{strong:thm:strong} give the asserted instability and
nonrepresentability.\qedhere
\end{enumerate}
\end{proof}

\begin{remark}[A prime-indexed family]
Replacing $2,4$ by $p,p^2$ gives the same construction for every prime
$p$: use $\rho_H(b)=b\otimes g+pa\otimes x$, $N=\Z/p^2$, and
$R=\Z/p$. The problematic summand is $N[p]b$, its generator has
cone image $pb\otimes v$, and the reduced comparison again has matrix
$\operatorname{diag}(1,1,1,0)$. The proofs above apply with these
substitutions.
\end{remark}

\section{Further applications: generalized multicategories and polycategories}\label{sec:kleisli}

The bicomodule application required reversal because its primitive maps
have the opposite arity to multiplication. In an original virtual double
category the monoid cells already have the required arity. Conservativity
therefore gives a direct comparison with generalized multicategories.
The distributor application likewise keeps the forward variance of its
module multicells when applying the completion for the second time.
This section records that comparison and the separate example supplied by
Garner's two-sided Kleisli construction.

\begin{definition}[Raw global forest calculus]
For any virtual double category $\mathbb K$, define $\Raw(\mathbb K)$
using paths as horizontal arrows. A square with $m>0$ outputs consists
of a subdivision of its upper path into $m$ consecutive, possibly empty
blocks, a vertical arrow at each cut vertex, and $m$ original cells
whose lateral arrows are those successive cut arrows. A square with no
outputs is a structural empty square: its upper path is empty and its
two exterior maps are the same vertical arrow. Composition substitutes
forests vertically and concatenates them horizontally.
\end{definition}

\begin{lemma}\label{lem:raw}
The raw forest calculus is a strict double category with paths as
horizontal arrows, and
\[
 \Raw(\mathbb K)_{f,g}(\Gamma;H)=\mathbb K_{f,g}(\Gamma;H).
\]
It therefore preserves and reflects formal categories and their
homomorphisms on atomic carriers.
\end{lemma}
\begin{proof}\leavevmode
\begin{enumerate}
\item \textbf{Verify substitution and interchange.} Substitution composes the arrows at cut vertices as well as the root
cells. Associativity is exactly virtual substitution associativity;
substitution on disjoint blocks commutes, proving interchange.

\item \textbf{Handle nullary and identity squares.} Nullary blocks use the degenerate virtual pastings, and the structural
empty square on $f$ is its horizontal identity.

\item \textbf{Recover the single-output fragment.} For one output there
is one block and its two cut arrows are the specified $f,g$, so the
displayed equality follows.

\item \textbf{Retain monoids and homomorphisms.} All monoid and homomorphism operations lie
in this fragment, which proves the last statement.\qedhere
\end{enumerate}
\end{proof}

\begin{theorem}[Kleisli recovery]\label{thm:kleisli}
Let $\mathbb K=H\text{-}\mathrm{Kl}(\mathbb X,T)$ be the horizontal
Kleisli virtual double category of a monad on a virtual double category.
There is an isomorphism
\[
 \FCat(\Up\Raw(\mathbb K))\cong T\text{-}\Mon(\mathbb X).
\]
When $\mathbb K$ is a virtual equipment, its refined realization
$\E_{\mathrm{poly}}(\mathbb K)$ gives the same comparison and also
preserves its represented composites and boundary transport.
\end{theorem}
\begin{proof}\leavevmode
\begin{enumerate}
\item \textbf{Identify the Kleisli monoid data.} A horizontal Kleisli arrow from $X$ to $Y$ is the original arrow
$X\harr TY$ with the Kleisli typing. By the definition of a $T$-monoid
\cite[Definition 4.3]{CruttwellShulman2010}, its multiplication and
unit are the monoid cells in $\mathbb K$.

\item \textbf{Pass to the path calculus and reduct.} Apply \cref{lem:raw} and
\cref{thm:reduct}.

\item \textbf{Apply equipment completion when available.} If $\mathbb K$ is an equipment, apply
\cref{thm:completion} instead; its inverted comparisons are precisely
the intrinsic ones by \cref{thm:detection}.\qedhere
\end{enumerate}
\end{proof}

\begin{proposition}[Normalization is tracked]\label{prop:normalization}
For a monad $T$ on a virtual equipment $\mathbb X$, the preceding
comparison restricts to normalized $T$-monoids. More explicitly, the
original unit square
$U_X\cell A$ with sides $(1_X,\eta_X)$ is cartesian in $\mathbb X$
if and only if its image is cartesian in $\E(\mathbb X)$.
\end{proposition}
\begin{proof}\leavevmode
\begin{enumerate}
\item \textbf{Identify the normalization test.} Normalization is the cartesianness of this unit square
\cite[Definition 8.2]{CruttwellShulman2010}. Its top is atomic and its
bottom is atomic in the original equipment $\mathbb X$.

\item \textbf{Preserve and reflect cartesianness.} \Cref{thm:globalembedding} preserves and reflects all the single-output
factorization tests defining cartesianness.

\item \textbf{Retain the original unit square.} The object correspondence in
\cref{thm:kleisli} retains the original unit square as data and
therefore retains this property.\qedhere
\end{enumerate}
\end{proof}

Here cartesianness is tested in $\mathbb X$, as in the definition of
normalization. No extension of $T$ across a localization is used to make
this test. This matters when $T$ itself does not satisfy the extension
criterion of \cref{thm:functoriality}.

\begin{proposition}[The symmetric-polycategory benchmark]
Garner's description of symmetric polycategories as monads on discrete
objects in a two-sided Kleisli bicategory is a further instance of
formal-category data in a poly-virtual calculus.
\end{proposition}
\begin{proof}\leavevmode
\begin{enumerate}
\item \textbf{Identify the Kleisli bicategory and its monads.} Garner constructs the two-sided Kleisli bicategory using a
pseudo-distributive law whose suitable matchings have associated graphs
that are connected, acyclic, and have no multiple edges
\cite[Proposition 9 and Definition 14]{Garner2008}, and identifies its
monads on discrete objects with symmetric polycategories
\cite[Theorem 42]{Garner2008}.

\item \textbf{Construct its coherent path presentation.} For any bicategory, fix a bracketing of each finite path of $1$-arrows
and take cells between the corresponding bracketed composites. Rectangular
evaluation is composition followed by the canonical associator and unitor
comparisons; bicategorical coherence makes the result independent of the
evaluation order \cite[Section 1.5]{Leinster2004}.

\item \textbf{Identify monads in the reduct.} A monad on
one original $1$-arrow is exactly \cref{eq:shape,eq:monadlaws} in
this path presentation. Apply \cref{thm:reduct}.

\item \textbf{Retain the given symmetric structure.} The symmetric
permutations here are part of Garner's Kleisli data; this argument does
not insert exchanges into our planar wire calculus.\qedhere
\end{enumerate}
\end{proof}

\section{Universal characterization and comparison with other envelopes}\label{sec:comparison}

\begin{theorem}[Characterization of the construction]
The globular merge completion is characterized by the following
presentation: freely form forests of the entire typed single-output
calculus and make every member of $S(P)$, and hence each of its horizontal
contextual instances, reversible.
It adds no original single-output cells and identifies no such cells.
Its poly-virtual reduct retains all cells and equalities, the tree
operations, and coherent boundary transport; it retains the entire
formal-category and bimodule-action theory.
Applied to the input $\Ddist(\V)$, the same characterization gives
$\FDist(\V)$ and preserves all primitive distributor multicells,
including their changing-base instances.
\end{theorem}
\begin{proof}\leavevmode
\begin{enumerate}
\item \textbf{Identify the presentation and reversible comparisons.} The presentation and its uniqueness are \cref{thm:localUP}.
The exact recognition of reversible atomic comparisons is
\cref{thm:detection}.

\item \textbf{Apply both conservativity theorems.} The two conservativity assertions are
\cref{thm:conservative,thm:globalembedding}.

\item \textbf{Pass to the reduct.} For the reduct use
\cref{thm:reduct,prop:mates,thm:bimodules}.

\item \textbf{Apply the distributor instance.} For the distributor instance apply \cref{thm:distcompletion}.\qedhere
\end{enumerate}
\end{proof}

The path presentation in \cref{prop:pathpresentation} is the merge
form of the two-sided calculus; compare the distinction between virtual
and implicit boundary operations in
\cite[Remark 1.1]{FairbanksShulman2025}. An independently defined module
or presheaf envelope requires the following comparison.

\begin{proposition}[A concrete comparison test]\label{prop:envelopetest}
Let $J:P\to\mathcal M$ be a globular interpretation in a path
presentation of another envelope. If $J$ sends every $s\in S(P)$
to an invertible comparison, there is a unique induced interpretation
$\bar J:L(P)\to\mathcal M$. It is an equivalence on a hom-category
precisely when it is fully faithful there and every object of that
hom-category is isomorphic to an interpreted path.
\end{proposition}
\begin{proof}\leavevmode
\begin{enumerate}
\item \textbf{Obtain the extension.} The extension statement is \cref{thm:localUP}.

\item \textbf{Construct a quasi-inverse.} A fully faithful
and essentially surjective functor is an equivalence: choose, for
each target object, an isomorphism from an object in the image; use
full faithfulness to pull back the conjugated morphisms. The resulting
functor is a quasi-inverse by those same isomorphisms.

\item \textbf{Prove necessity and apply the criterion.} Conversely an
equivalence is fully faithful and essentially surjective. Applied
to $\bar J$ this gives the stated test.\qedhere
\end{enumerate}
\end{proof}

An identification with a module or Yoneda envelope therefore requires
an interpretation that satisfies this test. The notation $\Mod$ used
here refers throughout to the monoids-and-bimodules construction. It
first supplies the balanced bicomodule calculus from $\V^\op$, and
then supplies distributor multicells from the path-carrier equipment
$\Kcal(\V)$. These two uses are explicit and do not identify the
result with an independently defined presheaf envelope. In particular,
\cref{rem:secondcomparison} explains why forgetting distributor actions
need not extend across the second localization.

The next concrete comparisons suggested by these results are to construct
such an interpretation for a chosen presheaf or module envelope, to study
the representations of distributor paths beyond the coequalizer
construction of \cref{thm:distcoequalizer}, and to replace the planar
tree shapes by other specified arity systems while retaining an analogue
of backward evaluation. An intrinsic distributor construction starting
only from an abstract tree calculus is another question: the construction
proved here uses the merge realization already provided for $\Q(\V)$.

\section{Conclusion}\label{sec:conclusion}

Finite bicomodule paths specify category and action boundaries in every
monoidal category. Stable cotensors represent those boundaries and recover
the usual operations, equations, and changing-base functors. The two
conservative completions retain the original balanced bicomodule and
distributor cells while making intrinsic contextual representations
reversible. Cartesian factorization supplies distributor restrictions;
contextually preserved action coequalizers represent their composition.

The examples separate two failures. Finite group-like bases in $\Ab$
have stable cotensors and recover the enriched equipment even though
$\Ab$ is not regular. The triangular example has stable category
composition but an action boundary with no contextual atomic
representative. Its nonzero action is constructed by restriction along
a formal functor on the same actual bicomodule. Thus the formal calculus
also supports operations beyond the stably represented sector.

\appendix
\section{Pasting conventions and admissible elementary shapes}\label{app:shapes}

The definition using rectangular schemes is finite and determines every
admissible operation. The elementary globular cases can also be listed
as follows; each formula requires matching vertex labels.
\begin{enumerate}
\item Input substitution:
$\alpha:\Gamma\cell M$, $\beta:(L,M,R)\cell\Delta$ give
$(L,\Gamma,R)\cell\Delta$.
\item Output substitution:
$\alpha:\Gamma\cell(L,M,R)$, $\beta:M\cell\Delta$ give
$\Gamma\cell(L,\Delta,R)$.
\item A right corner cut is \cref{eq:cornercut}.
\item A left corner cut has premises
$\Gamma\cell(M,\Delta)$ and $(\Lambda,M)\cell\Sigma$ and
conclusion $(\Lambda,\Gamma)\cell(\Sigma,\Delta)$.
\end{enumerate}
These cases, with their overlaps, are exactly
\cref{eq:cutplacement}. Arbitrary vertical labels are governed by the
rectangular matching rule, not by suppressing the labels in this list.
For example, to substitute $n$ cells above a virtual cell one specifies
all $n+1$ interface vertical arrows together; adjacent members must use
the very same map at their common vertex.

The equations defining evaluation have an equally finite structural
description: deleting structural strips does not change evaluation;
evaluating a subtree and replacing it by its value does not change
evaluation; and equivalent rectangular subdivisions have the same
value. Associativity of nested substitutions is the second rule.
For disjoint subtrees it applies twice in either order. All higher
instances follow by induction on the number of rectangles. In
particular, no equation asserts equality of differently typed
pastings or supplies a cut where an exterior directed boundary fails
to exist.

\section{Explicit adjunction calculations}\label{app:bindings}

This appendix records the algebra needed in the global construction.
For composable $f,u,w$, the companion comparisons satisfy
\begin{equation}\label{eq:kassoc}
 k_{uf,w}(k_{f,u}\wh1_{w_*})
   =k_{f,wu}(1_{f_*}\wh k_{u,w}).
\end{equation}
The conjoint equation is its reversed-order counterpart. Both sides
of \cref{eq:kassoc} carry the pasted binding pair to the binding
pair of $wuf$; the uniqueness argument in \cref{lem:binding}
therefore proves the equation. Applied to four factors, this same
characterization proves every association comparison equal, without
an additional coherence choice.

For a globular $a:\Gamma\cell f_*\Delta g^*$ write
\[
 \mathfrak m(a)=(e_f\wh1_\Delta\wh e_g)
                   (1_{f^*}\wh a\wh1_{g_*}).
\]
For $b:f^*\Gamma g_*\cell\Delta$, its inverse mate is
\[
 \mathfrak n(b)=
 (1_{f_*}\wh b\wh1_{g^*})(j_f\wh1_\Gamma\wh j_g).
\]
Expanding $\mathfrak m\mathfrak n(b)$ makes the $f$ terms contract
by $(e_f\wh1_{f^*})(1_{f^*}\wh j_f)=1_{f^*}$ and the $g$
terms by $(1_{g_*}\wh e_g)(j_g\wh1_{g_*})=1_{g_*}$.
The result is $b$. Expanding $\mathfrak n\mathfrak m(a)$ uses the
other two triangle identities and gives $a$. These expansions also
prove naturality: a whiskered cell can be moved past a contraction
on a disjoint factor by interchange.

For the interchange calculation in \cref{thm:globalcoherence}, the
only nontrivial local equality is
\[
 e_{vg}(l_{g,v}\wh k_{g,v})
     =e_v(1_{v^*}\wh e_g\wh1_{v_*})
 :v^*g^*g_*v_*\cell\emp{\operatorname{cod}(v)}.
\]
This is \cref{eq:counitcomp} with its invertible comparisons moved
to the other side. It covers every interior seam of a rectangular
grid. At an empty horizontal boundary the corresponding equality is
\cref{eq:unitcomp}. These two identities, the triangles, and
\cref{eq:kassoc} account for all identity, association, and
interchange checks in the realization.

\section{Finite terms and size}\label{app:size}

\begin{proposition}
For a small input signature and small original cell collections, the
forest envelope, localization, and global realization have small cell
collections. For a class-sized input the same specifications define a
finite-term calculus in an encompassing universe; they do not by
themselves establish local smallness in the original universe.
The same conclusions apply to the path-carrier, distributor, and second
completion constructions.
\end{proposition}
\begin{proof}\leavevmode
\begin{enumerate}
\item \textbf{Bound the forest collections.} For small data, finite paths and finite subdivisions form sets.
Forests are finite products of original cell sets, indexed by a set
of subdivisions, hence form a set.

\item \textbf{Bound the localization.} The collection $S(P)$ is a
subcollection of the original cells, so is a set as well. Finite
terms in these generators and their formal inverses form a set;
the congruence generated by the stated relations is an equivalence
relation on it, and its quotient is a set.

\item \textbf{Bound the global cell collections.} For a fixed vertical arrow $f$,
the binding presentations in $\mathcal B(f)$ form a subcollection of a
finite product of the original horizontal-arrow and cell sets, and hence a
set. Therefore the coproduct and quotient in \cref{eq:choicefreecells} are
sets. After a local binding presentation is fixed, each global cell
collection is identified with the localized hom-set in
\cref{eq:globalcells}.

\item \textbf{Bound the distributor completion.} For the distributor construction on small data, paths and the cells of
the carrier equipment form sets. Formal categories and distributors are
specified by finitely many carriers and cells satisfying equations, so
they form subcollections of sets of such tuples. Distributor multicells
are subsets of the corresponding carrier cell sets. The forest and
localization argument therefore applies a second time. Restricting on
the original atomic monoid objects does not change this conclusion.

\item \textbf{Treat large data in the larger universe.} For large data the same description uses finite terms labelled in
the larger universe. Although every individual cell has a finite
description, arbitrarily many intermediate atomic arrows can occur
in descriptions with fixed many-sided endpoints. Finite length
alone therefore does not give a small collection in the smaller
universe.

\item \textbf{Identify the fragments with preserved size.} The conservativity theorem does give exactly the original
size for the single-output fragment and, after reversal, for the
single-input fragment.
The second completion similarly preserves exactly the size of the
primitive single-output distributor collections, without deducing local
smallness of its general many-sided profiles in a smaller universe.\qedhere
\end{enumerate}
\end{proof}

\begin{thebibliography}{Had19}

\bibitem[Agu97]{Aguiar1997}
Marcelo Aguiar,
\textbf{Internal categories and quantum groups},
Ph.D. thesis, Cornell University, 1997.
\url{https://pi.math.cornell.edu/~maguiar/thesis2.pdf}.

\bibitem[CS10]{CruttwellShulman2010}
G.~S.~H. Cruttwell and Michael~A. Shulman,
\textbf{A unified framework for generalized multicategories},
Theory and Applications of Categories \textbf{24} (2010), no.~21,
580--655.
\href{https://arxiv.org/abs/0907.2460}{arXiv:0907.2460}.

\bibitem[FS26]{FairbanksShulman2025}
Aaron David Fairbanks and Michael Shulman,
\textbf{Doubly weak double categories},
Applied Categorical Structures \textbf{34} (2026), article 26.
\href{https://doi.org/10.1007/s10485-026-09863-1}
{doi:10.1007/s10485-026-09863-1}.

\bibitem[Gar08]{Garner2008}
Richard Garner,
\textbf{Polycategories via pseudo-distributive laws},
Advances in Mathematics \textbf{218} (2008), 781--827.
\href{https://arxiv.org/abs/math/0606735}{arXiv:math/0606735}.

\bibitem[Had19]{Hadzihasanovic2019}
Amar Hadzihasanovic,
\textbf{Weak units, universal cells, and coherence via universality for
bicategories},
Theory and Applications of Categories \textbf{34} (2019), no.~29,
883--960.
\href{https://arxiv.org/abs/1803.06086}{arXiv:1803.06086}.

\bibitem[Lei04]{Leinster2004}
Tom Leinster,
\textbf{Higher operads, higher categories},
London Mathematical Society Lecture Note Series, vol.~298,
Cambridge University Press, Cambridge, 2004.
\href{https://doi.org/10.1017/CBO9780511525896}
{doi:10.1017/CBO9780511525896}.

\bibitem[Shu08]{Shulman2008}
Michael Shulman,
\textbf{Framed bicategories and monoidal fibrations},
Theory and Applications of Categories \textbf{20} (2008), no.~18,
650--738.
\href{https://arxiv.org/abs/0706.1286}{arXiv:0706.1286}.

\end{thebibliography}
\end{document}
